\documentclass[11pt,oneside]{article}
\usepackage[T1]{fontenc}
\usepackage[utf8]{inputenc}
\usepackage{lmodern}
\usepackage{microtype}
\usepackage[a4paper,margin=30mm,headheight=14pt]{geometry}
\usepackage{amsmath,amssymb,amsthm,mathtools}
\usepackage{booktabs,longtable,tabularx,array,multirow}
\usepackage{enumitem}
\usepackage{xcolor}
\usepackage{listings}
\usepackage{fancyhdr}
\usepackage{graphicx}
\usepackage{url}
\usepackage[hidelinks,bookmarksnumbered]{hyperref}
\usepackage[nameinlink,noabbrev]{cleveref}
\usepackage{caption}
\usepackage{etoolbox}

\definecolor{codebg}{RGB}{247,247,247}
\definecolor{darkblue}{RGB}{25,58,112}
\definecolor{darkgreen}{RGB}{35,100,55}
\definecolor{darkred}{RGB}{135,45,45}
\hypersetup{colorlinks=true,linkcolor=darkblue,citecolor=darkgreen,urlcolor=darkred}

\lstset{
  basicstyle=\ttfamily\small,
  backgroundcolor=\color{codebg},
  frame=single,
  framerule=0.2pt,
  breaklines=true,
  columns=fullflexible,
  keepspaces=true,
  showstringspaces=false,
  xleftmargin=0.5em,
  xrightmargin=0.5em
}

\newtheorem{theorem}{Theorem}[section]
\newtheorem{lemma}[theorem]{Lemma}
\newtheorem{proposition}[theorem]{Proposition}
\newtheorem{corollary}[theorem]{Corollary}
\newtheorem{principle}[theorem]{Principle}
\newtheorem{barrier}[theorem]{Barrier}
\newtheorem{schema}[theorem]{Schema}
\theoremstyle{definition}
\newtheorem{definition}[theorem]{Definition}
\newtheorem{example}[theorem]{Example}
\newtheorem{construction}[theorem]{Construction}
\newtheorem{interface}[theorem]{Interface}
\theoremstyle{remark}
\newtheorem{remark}[theorem]{Remark}
\newtheorem{status}[theorem]{Status}

\newcommand{\N}{\mathbb N}
\newcommand{\Z}{\mathbb Z}
\newcommand{\Q}{\mathbb Q}
\newcommand{\B}{\{0,1\}}
\newcommand{\iPRA}{\ensuremath{\mathrm{iPRA}}}
\newcommand{\FRNF}{\ensuremath{\mathrm{FRNF}}}
\newcommand{\PRAAx}{\ensuremath{\mathsf{PRAAx}}}
\newcommand{\Prf}{\ensuremath{\mathsf{Prf}}}
\newcommand{\PTm}{\ensuremath{\mathsf{PTm}}}
\newcommand{\Fm}{\ensuremath{\mathsf{Fm}}}
\newcommand{\QF}{\ensuremath{\mathsf{QF}}}
\newcommand{\code}[1]{\texttt{\detokenize{#1}}}
\newcommand{\abs}[1]{\lvert #1\rvert}
\newcommand{\norm}[1]{\lVert #1\rVert}
\newcommand{\diam}{\operatorname{diam}}
\newcommand{\sg}{\operatorname{sg}}
\newcommand{\least}{\operatorname{least}}
\newcommand{\prefixmax}{\operatorname{pmax}}
\newcommand{\Valid}{\operatorname{Valid}}
\newcommand{\Read}{\operatorname{Read}}
\newcommand{\Just}{\operatorname{Just}}
\newcommand{\Viable}{\operatorname{Viable}}
\newcommand{\Obs}{\operatorname{Obs}}
\newcommand{\rank}{\operatorname{rk}}
\newcommand{\Realize}{\operatorname{Realize}}
\newcommand{\Eval}{\operatorname{Eval}}
\newcommand{\Bool}{\operatorname{Bool}}
\newcommand{\eps}{\varepsilon}

\pagestyle{fancy}
\fancyhf{}
\fancyhead[L]{Finitary Realizer Normal Form for iPRA}
\fancyhead[R]{\thepage}
\fancyfoot[C]{Research package, 16 November 2026}

\title{\textbf{Finitary Realizer Normal Form}\\
A proof-producing framework for embedding mathematics into\\
intuitionistic primitive recursive arithmetic}
\author{Research report prepared for the ProveIt iPRA development}
\date{16 November 2026}

\begin{document}
\maketitle

\begin{abstract}
The ProveIt repository contains a growing collection of finitary reconstructions of
classical theorems, together with a checked core of ordinary intuitionistic primitive
recursive arithmetic and many reports identifying recurring obstructions.  The
collection already exhibits a remarkably stable pattern: exact finite arithmetic is
straightforward, while the difficult passages concern represented infinite objects,
adaptive precision, nonunique witnesses, compactness, branch selection, and the gap
between a host-language theorem and an object-language derivation.

This report isolates that pattern as a common intermediate language, the
\emph{Finitary Realizer Normal Form} (\FRNF).  An \FRNF\ proof is a fixed finite graph
of statically specialized primitive-recursive transducers, bounded searches, clocked
loops with quantifier-free invariants, finite relational witness clouds, stable
observers, and explicit producer boundaries.  Every edge carries two primitive-
recursive moduli: a \emph{computation horizon}, recording which input rows are read,
and a usually larger \emph{justification horizon}, recording which input contracts
are required to prove the output correct.  A reverse demand pass is run before any
induction.  The resulting finite table is Booleanized, one clocked trace is verified
by quantifier-free induction, and the local embedded derivations are composed by
simultaneous cut into a term of type \(\Prf\,\PRAAx\).  Nonunique approximate
witnesses are not prematurely forced into exact selectors: backward viability gives
finite coherent paths, stable observables erase internal choices, and exact branch
labels are recovered only from explicit separation data.

The report gives complete human proofs of bounded-formula Booleanization, least
bounded search, dual-horizon composition, reverse demand closure, trace
internalization, backward viability, shrinking finite-horizon coherence, majorized
descent, finite-net minimization, and guarded inversion.  It proves two representation
barriers---there is no total primitive-recursive universal evaluator, and a uniform
cardinality-to-cutoff operator would decide the halting problem---and uses them to
separate genuine mathematical inputs from mere formalization work.  This resolves a
large class of reported difficulties and gives exact conditional endpoints for the
remaining ones, including arbitrary Cauchy Puiseux data and general Siegel/Faltings
height bounds.

The accompanying artifacts include exact-arithmetic Python reference algorithms with
eleven regression tests, a repository-aligned Lean module containing the logical
\(\Prf\,\PRAAx\) combinators, explicit Lean targets for the remaining arithmetic
compiler, machine-readable theorem and obstacle ledgers, and this source.  The Lean
files have not been compiled in the artifact environment; visible placeholders are
retained precisely where object-level implementation remains.
\end{abstract}

\vspace{0.8em}
\noindent\textbf{Repository snapshot.}
The audit was made against the public ProveIt repository on 16 November 2026, with
the current PrimitiveRecursiveArithmetic README at commit
\code{67f5216e0b1337b533b43064262aeff268bc56db}.  Paths and declaration names in this
report refer to that snapshot unless stated otherwise.

\tableofcontents
\clearpage

\section{Scope, contribution, and claim discipline}
\label{sec:scope}

\subsection{What problem is being solved?}

The intended target is not the informal statement that ``ordinary mathematics can be
encoded in arithmetic.''  Such a statement is too coarse to guide a formalization.  A
usable framework must answer five concrete questions for every theorem.

\begin{enumerate}[label=\textbf{Q\arabic*.},leftmargin=3.4em]
\item What finite numerical data represent each mathematical object at the precision
actually consumed by the proof?
\item Which primitive-recursive programs compute the finite outputs, and how are their
adaptive queries bounded before induction begins?
\item Which formulas express correctness in the first-order language of \(\iPRA\), and
which of those formulas are quantifier-free induction invariants?
\item How are local correctness derivations composed into an actual object-theory
proof, rather than merely a theorem about the intended standard model?
\item Which pieces of data are generated by a primitive-recursive producer, and which
remain explicit mathematical assumptions such as a modulus, a height bound, or a
separation constant?
\end{enumerate}

The existing ProveIt packets repeatedly answer subsets of these questions \cite{proveit}.  The
checked \code{Derivations} library shows that substantial exact mathematics can be
carried all the way to terms of type
\[
   ([] : \mathrm{List}(\Fm\,\PTm\,0))\ \vdash_{\PRAAx}\ \varphi.
\]
The unverified packet corpus contains many sophisticated quantitative proofs and
several overlapping general frameworks.  The bottleneck is therefore not a shortage
of high-level ideas.  It is the absence of one normal form that is simultaneously
mathematically expressive, honest about representation barriers, and designed to emit
embedded derivations.

\subsection{The proposed normal form}

We call the normal form developed here \emph{Finitary Realizer Normal Form}, abbreviated
\(\FRNF\).  Its central design decision is that a proof is compiled twice.

First, at the \emph{semantic compilation} level, the proof is specialized to the fixed
primitive-recursive routines that occur in the theorem.  Higher-type language is
eliminated in favor of explicit finite observations, moduli, bounds, and relations.
All adaptive behavior is defunctionalized into a finite state machine with an explicit
clock.  The compiler computes a finite closure of every row of every input name that
is needed either for computation or for justification.

Second, at the \emph{object-proof compilation} level, the closed demand table is
represented by primitive-recursive symbols and first-order formulas.  Bounded
conditions are replaced by characteristic terms.  A single quantifier-free invariant
tracks the clocked trace.  Local derivations are then assembled in topological order.
The final output is not a Boolean test, a Python success message, or a Lean theorem
about natural numbers.  It is an inhabitant of \(\Prf\,\PRAAx\) for the intended
formula.

\subsection{What is new in the synthesis?}

Several ingredients have precedents in the repository and in proof theory.  Bounded
search, primitive-recursive coding, proof interpretations, moduli, and finite nets are
standard tools.  The contribution here is their organization around four requirements
that are easy to state but rarely enforced together.

\begin{enumerate}[label=(\roman*),leftmargin=2.4em]
\item \textbf{Dual horizons.}  Every interface distinguishes the input rows needed to
compute an answer from the input rows needed to prove the answer valid.  Composition
propagates the second quantity, not merely the first.
\item \textbf{Relational intermediate values.}  Approximation, branch ambiguity, and
collisions are represented by finite candidate relations.  Deterministic selection is
introduced only when the theorem proves that the observable is choice-independent or
an explicit gap isolates a branch.
\item \textbf{A formula firewall.}  Host-language semantics and object-language
provability are separate fields of every interface.  There is no coercion from a Lean
proposition to \(\Prf\,\PRAAx\).
\item \textbf{Boundary preservation.}  A finite verifier is never advertised as a
producer of a height bound, modulus, zero-free region, or exact algebraic splitting.
Such data remain visible hypotheses until a primitive-recursive producer and its
object proof have themselves been supplied.
\end{enumerate}

The dual-horizon requirement is especially important.  In a represented proof, the
program may read a row \(a(r(d))\) to produce output precision \(d\), while the proof
that this row is correct may require coherence comparisons at \(a(j(d))\) with
\(j(d)>r(d)\).  Collapsing these into one ``use'' bound is safe only if the bound is
defined to dominate both.  Many informal formalization plans silently track the first
and need the second.

\subsection{Claim classes}

To keep the package auditable, every result belongs to one of the following classes.

\begin{description}[style=nextline,leftmargin=2em]
\item[Mathematical theorem.] A complete ordinary proof is given in this report.
\item[Executable finite check.] The standard-library Python reference implementation
runs and its regression test passes.  This validates the finite algorithm on the test
cases, not its object-theory proof.
\item[Host-checked Lean theorem.] Lean accepts a theorem in its own metatheory.  No such
claim is made here because Lean was not installed in the artifact runtime.
\item[Object-checked theorem.] Lean accepts a declaration whose conclusion is an
embedded \(\Prf\,\PRAAx\) derivation, and an axiom audit has been inspected.  This
status belongs to the upstream \code{Derivations} results that we reuse; it is not
claimed for the new generic compiler.
\item[Blueprint target.] The exact statement, definitions, dependencies, and proof
route are supplied, but placeholders remain.
\item[Boundary or obstruction.] The framework proves that an input cannot be generated
by the proposed generic mechanism, or that extra data are logically required.
\end{description}

This discipline matters because an executable checker and an object proof answer
different questions.  The former asks whether a concrete finite certificate passes a
program.  The latter asks whether the correctness of that program has been derived
using only the axioms and induction allowed by \(\iPRA\).

\subsection{Relation to proof interpretations}

The guiding philosophy is compatible with G\"odel's functional interpretation and its
modern proof-mining descendants: apparently infinitary assertions are transformed into
finite computational data and quantitative bounds \cite{goedel1958,kohlenbach2008}.
The present goal is narrower and more syntactic.  We do not provide a global
interpretation of a strong ambient theory.  Instead, we give a compositional target
language for hand-developed finitary proofs and a route from that language to the
particular embedded calculus used by ProveIt.  One may view \(\FRNF\) as a proof-
engineering normal form informed by proof interpretation rather than as a new
conservation theorem.

\subsection{The main theorem in informal form}

The report culminates in the following schema, made precise in
\cref{sec:compilation-theorem}.

\begin{schema}[FRNF compilation, informal]
Fix a finite acyclic graph whose nodes are statically specialized primitive-recursive
transducers, bounded searches, clocked loops, finite relational witnesses, stable
observers, or explicit boundary assumptions.  Suppose every edge has primitive-
recursive computation and justification horizons, every loop has a quantifier-free
invariant with embedded base and step derivations, and every node has an embedded local
correctness derivation from its declared predecessors.  Then for every fixed final
demand the graph compiles to an embedded \(\iPRA\) derivation of the final contract.
The only mathematical assumptions in the result are the boundary formulas explicitly
listed by the graph.
\end{schema}

The theorem is deliberately a schema over fixed syntax.  It does not quantify inside
\(\iPRA\) over arbitrary primitive-recursive program codes.  \Cref{sec:static}
proves why that restriction is essential rather than cosmetic.

\section{The target theory and the standard of success}
\label{sec:ipra-target}

\subsection{Working presentation of iPRA}

For this report, \(\iPRA\) is first-order intuitionistic arithmetic with equality, a
function symbol for every primitive-recursive function and its defining equations,
successor axioms, and induction restricted to quantifier-free formulas.  This is the
presentation implemented by the ProveIt type \(\Prf\,\PRAAx\).  The exact choice of a
primitive-recursive basis is immaterial up to definitional extension, but it is not
immaterial to a formalization: every algorithm mentioned in a proof must eventually
be represented by a fixed symbol or a term assembled from such symbols.

The restriction to quantifier-free induction does not make the theory incapable of
proving quantified theorems.  It constrains the induction invariant.  A typical proof
first encodes a finite computation, proves a quantifier-free invariant for the encoded
state, specializes the invariant at a primitive-recursive clock, and only then uses
ordinary intuitionistic introduction and elimination rules to derive the final
quantified statement.

The repository supplies two complementary targets.

\begin{itemize}
\item \code{ConciseIPRA} formalizes ordinary \(\iPRA\), a very small native scalar
calculus, and translations between them.  Its theorem \code{pra_iff_native} exports an
ordinary embedded derivation to the native core.
\item \code{IPRADerivations} develops arithmetic directly in
\(\Prf\,\PRAAx\): semiring identities, order, truncated subtraction, pairing,
simultaneous recursion, bounded sums and search, division, gcd, B\'ezout, descent, and
other reusable components.
\end{itemize}

The framework below is therefore not a replacement foundation.  It is a compiler
front end whose output is designed to enter this checked library.

\subsection{Object theory versus metatheory}

We use three levels of assertion.  Confusing them is the single most damaging source
of false completion claims.

\paragraph{Level 0: executable mathematics.}
A primitive-recursive algorithm is described extensionally and may be run in Python,
Lean, or another language.  One can test examples and inspect traces.  This level is
useful for finding off-by-one errors and validating finite combinatorics.

\paragraph{Level 1: host semantics.}
Lean proves a theorem such as
\[
   \forall x:\N,\quad \operatorname{run}(x)\ \text{satisfies}\ \operatorname{Spec}(x).
\]
This theorem lives in Lean's dependent type theory and may use host induction,
quotients, classical choice, or a rich library.  It is evidence that the mathematics
has been formalized, but it is not automatically a theorem of \(\iPRA\).

\paragraph{Level 2: embedded object proof.}
Lean constructs a term of a type such as
\[
   ([] : \mathrm{List}(\Fm\,\PTm\,0))\vdash_{\PRAAx}\varphi.
\]
The inductive structure of this term is the formal derivation.  Mathlib may be used in
the metatheory to build and check it, but the object proof contains only instances of
\(\PRAAx\) and the deduction rules.  This is the endpoint required here.

\begin{principle}[No semantic promotion]
A host theorem about the standard natural numbers may justify the design of a
primitive-recursive symbol, but it may not be used as an object theorem unless a
translation theorem explicitly constructs the corresponding \(\Prf\,\PRAAx\)
derivation.
\end{principle}

This principle does not prohibit semantic reasoning.  It assigns it the correct role.
For example, a host proof may establish that a proposed rank really decreases; the
formalization must still encode the rank and derive the decrease formula in the object
theory before using quantifier-free induction.

\subsection{The shape of acceptable induction}

Suppose a bounded process has state \(s_t\) at time \(t\), encoded by a number, and a
fixed final demand \(d\).  An unacceptable induction hypothesis often has the form
\[
  \forall k\;\forall n\ge M(k),\quad \abs{a_n-a_{n+1}}\le 2^{-k}.
\]
It quantifies over a represented infinite object and is not quantifier-free.

The \(\FRNF\) replacement first computes a concrete horizon \(H(d)\).  It stores the
finite table
\[
  a(0),a(1),\ldots,a(H(d))
\]
and the finite collection of required inequalities as encoded numerical data.  The
invariant then states only equalities and inequalities among coded table entries,
clock values, and accumulator fields.  Universal quantifiers bounded by table length
are Booleanized into finite products or sums.  After the loop is closed, the proof
extracts the desired formula at demand \(d\).

The important move is not merely to say ``fix \(d\).''  One must close all dependencies
of \(d\) first.  If proving row \(d\) needs row \(d+3\), and proving that row needs an
anchor at \(2d+7\), the invariant must be built after those demands have been
propagated.  This is the purpose of dual horizons and reverse demand closure.

\subsection{Uniformity and schemas}

There are two legitimate forms of uniformity.

\begin{enumerate}[label=(\alph*)]
\item A single primitive-recursive symbol may accept numerical parameters such as a
matrix size, polynomial degree, requested precision, or finite graph code, provided
its algorithm is itself fixed.
\item A metatheoretic schema may associate a different primitive-recursive symbol and
derivation to each externally fixed program or theorem instance.
\end{enumerate}

What is not generally available is a total primitive-recursive evaluator
\(U(e,x)\) whose numerical code \(e\) ranges over all primitive-recursive programs and
which returns the value of every such program.  The diagonal argument in
\cref{sec:static} rules this out.  Thus, ``uniform over the degree'' can be a genuine
single theorem, while ``uniform over an arbitrary code of a primitive-recursive
coefficient evaluator'' must be interpreted as a static metatheoretic schema or be
restricted to a fixed, bounded interpreter.

\subsection{Constructive content of familiar mathematical language}

The following translations recur throughout the report.

\begin{center}
\begin{tabularx}{\textwidth}{@{}p{0.27\textwidth}X@{}}
\toprule
Informal phrase & Required \(\FRNF\) data \\
\midrule
``Let \(x\) be a real number'' & A fixed evaluator for rational approximants and a
coherence modulus, or a finite relational presentation. \\
``Choose a root'' & A finite root cloud plus residual and coverage certificates;
coherent choice only from motion/diameter data; exact labeling only from separation. \\
``By compactness'' & A primitive-recursive finite net, a coverage proof, and the
modulus converting net radius into output error. \\
``Since \(x\ne0\), divide by \(x\)'' & A positive rational apartness witness
\(\abs{x}\ge\lambda>0\). \\
``There are finitely many solutions'' & A height/cutoff producer or a complete-list
certificate, not just a cardinality bound. \\
``Run the algorithm until it stops'' & A primitive-recursive rank, clock, or other
majorant, with a QF decrease invariant and terminal freezing. \\
``The result is independent of choices'' & A stable-observer derivation comparing all
finite candidates, or a common invariant of the relation. \\
\bottomrule
\end{tabularx}
\end{center}

These are not bureaucratic substitutions.  They identify the exact data consumed by
constructive proofs and make hidden uses of choice, compactness, or excluded middle
visible.

\section{The formula firewall and the four-part interface}
\label{sec:firewall}

\subsection{Why an interface needs four parts}

A reusable component is often described by only a program and a semantic theorem.
That is sufficient in an ordinary programming library but insufficient for an
embedded weak theory.  We require four distinct fields.

\begin{definition}[FRNF component]
An \(\FRNF\) component from input presentation \(X\) to output presentation \(Y\)
consists of:
\begin{enumerate}[label=(\roman*),leftmargin=2.4em]
\item fixed primitive-recursive symbols computing finite output observations;
\item an object-language formula \(C_{X\to Y}\) expressing the component contract;
\item a constructor that, from embedded derivations of declared input contracts,
returns an embedded derivation of \(C_{X\to Y}\); and
\item an optional host-semantic theorem identifying the numerical contract with its
intended mathematical meaning.
\end{enumerate}
The first three fields are the \emph{kernel-facing interface}.  The fourth is the
\emph{semantic interpretation}.
\end{definition}

The separation between (iii) and (iv) is the formula firewall.  It may feel redundant
when the host theorem is elementary, but it allows components to be combined without
accidentally importing a strong metatheory into the object proof.

\subsection{Exact contracts}

The simplest component computes a term \(w(x)\) and proves a formula
\(Q(x,w(x))\) from a precondition \(P(x)\).  At the object level it has the shape
\[
    [P]\ \vdash_{\PRAAx}\ Q[w].
\]
Existential hiding then gives
\[
    [P]\ \vdash_{\PRAAx}\ \exists y\,Q(x,y),
\]
and implication introduction gives a closed contract.  Sequential composition is
substitution followed by cut.  The accompanying
\code{lean/FRNF/Kernel.lean} records exactly these logical patterns through a proof
schedule.

The crucial point is that an exact contract is not merely a structure containing a
Lean function \(w\) and a field of type \(Q(w)\).  Its witness is an object term
\(\PTm\), and its certificate is an embedded derivation.

\subsection{Relational contracts}

Many constructive constructions do not justify a single exact witness.  A polynomial
with a multiple root, an eigenvalue cluster, a chart overlap, or a discretized
optimization problem naturally yields a nonempty finite set of acceptable answers.
For such a component we use a bounded relation
\[
  R(x,i), \qquad i<B(x),
\]
with an object derivation of
\[
  P(x)\ \longrightarrow\ \exists i<B(x)\,R(x,i).
\]
The relation may carry a compatibility edge to the candidates at the next precision.
A downstream component must either be defined relationally on every valid candidate or
prove that its observable is stable across the candidate cloud.

This is stronger than choosing the least candidate immediately.  Least choice is
primitive recursive over a finite set, but it can destroy extensionality: two
presentations of the same represented object may enumerate candidates in different
orders.  Keeping the relation until the observation point allows one to prove
presentation-independent results.

\subsection{Conditional contracts and boundary nodes}

Suppose a finite algorithm enumerates all solutions below a supplied height bound
\(H(p)\).  Its correctness is a theorem of weak arithmetic once \(H\) is fixed and the
height theorem is given.  The statement should therefore be represented as
\[
  \operatorname{HeightBound}(p,H(p))
  \longrightarrow
  \operatorname{CompleteList}(p,\operatorname{enum}(p,H(p))).
\]
The height theorem is a boundary node.  It is not hidden in the implementation of
\code{enum}.  If a later project supplies an object proof of the height bound, ordinary
cut discharges the boundary.  Until then the endpoint remains a precise conditional
theorem rather than an incorrect claim of unconditional enumeration.

\begin{principle}[Producer/verifier separation]
A primitive-recursive verifier for a finite certificate does not become a
primitive-recursive producer of that certificate.  Production and verification are
separate nodes with separate proof obligations.
\end{principle}

This distinction is decisive in the general Siegel and Faltings packets, exceptional-
zero arguments in analytic number theory, arbitrary compactness arguments, and exact
splitting data for represented algebraic functions.

\subsection{Status fields as part of the mathematics}

A practical artifact should store, for every component, at least the following status
fields:

\begin{enumerate}[label=\arabic*.,leftmargin=2.2em]
\item semantic statement written;
\item finite reference implementation tested;
\item object formula defined;
\item local \(\Prf\,\PRAAx\) constructor implemented;
\item composed endpoint built;
\item native transfer built;
\item remaining boundary assumptions enumerated;
\item axiom audit inspected.
\end{enumerate}

These fields are not project-management decoration.  They express logically distinct
claims.  In particular, a compiled host theorem with a remaining \code{sorry} in the
object constructor is not a partial object proof; it is a complete result at one level
and an open result at another.

\subsection{Compositionality}

Let \(C_1,\ldots,C_m\) be component contracts arranged in topological order.  The local
proof for node \(i\) may use only the permanent hypotheses \(\Gamma\) and a declared
subset of earlier contracts.  If
\[
   \Delta_i,\Gamma\vdash_{\PRAAx} C_i
   \quad\text{and}\quad
   \Gamma\vdash_{\PRAAx}D\quad(D\in\Delta_i),
\]
then simultaneous cut yields \(\Gamma\vdash_{\PRAAx}C_i\).  Induction on the node list
therefore produces all graph contracts.  This simple observation is the logical core
of the proof schedule in \cref{sec:proof-schedule}; all substantive work lies in
making each local contract finite and QF-internalizable.

\section{Bounded formulas and Booleanization}
\label{sec:booleanization}

Quantifier-free induction is easiest to use when all finite side conditions are
represented by equations between primitive-recursive terms.  Bounded quantifiers do
not literally belong to the quantifier-free fragment, but every bounded formula has a
primitive-recursive characteristic function.  This familiar fact must be developed
constructively and in a form compatible with the embedded syntax.

\subsection{Bounded formula grammar}

Fix free variables \(\vec x\).  The bounded formulas are generated by
\[
\begin{split}
 A ::= {}& s=t \mid \bot \mid A\wedge A \mid A\vee A \mid A\to A \\
         &\mid \forall u<b(\vec x)\,A(u,\vec x)
          \mid \exists u<b(\vec x)\,A(u,\vec x),
\end{split}
\]
where all terms denote primitive-recursive functions.  Negation is \(A\to\bot\).
Because equality of naturals is decidable in \(\iPRA\), every such formula is
decidable, but we need more: a fixed characteristic term and embedded equivalence
proofs.

\subsection{Boolean arithmetic}

Let \([s=t]\in\B\) be the primitive-recursive equality test.  The checked
\code{Core/Pairing} module already supplies an equality term
\code{eqb} with object proofs for equal and unequal arguments.  Let
\(\sg(n)=0\) when \(n=0\) and \(1\) otherwise, and let
\(1\mathbin{\dot-}n\) denote truncated subtraction.

For Boolean inputs \(a,b\in\B\), define
\[
\begin{array}{rclcrcl}
 a\mathbin{\&}b&=&ab, &&
 a\mathbin{|}b&=&\sg(a+b),\\
 a\Rightarrow b&=&\sg((1\mathbin{\dot-}a)+b),&&
 \neg a&=&1\mathbin{\dot-}a.
\end{array}
\]
All four results are Boolean.  Intuitionistically,
\(a\Rightarrow b=1\) is equivalent to \(a=1\to b=1\): if \(a=0\), the
truncated difference contributes \(1\); if \(a=1\), the expression has sign
\(b\).  This finite case split is derivable from the successor axioms and
quantifier-free induction, and no excluded middle for arbitrary formulas is involved.

\subsection{Structural characteristic compiler}

\begin{construction}[Characteristic term]
Associate to every bounded formula \(A\) a primitive-recursive term
\(\chi_A(\vec x)\) by structural recursion:
\[
\begin{array}{rcl}
\chi_{s=t}&=&[s=t],\\
\chi_{\bot}&=&0,\\
\chi_{A\wedge B}&=&\chi_A\chi_B,\\
\chi_{A\vee B}&=&\sg(\chi_A+\chi_B),\\
\chi_{A\to B}&=&\sg((1\mathbin{\dot-}\chi_A)+\chi_B),\\
\chi_{\exists u<b A(u)}&=&\sg\!\left(\displaystyle\sum_{u<b}\chi_A(u)\right),\\
\chi_{\forall u<b A(u)}&=&\displaystyle\prod_{u<b}\chi_A(u).
\end{array}
\]
Empty sums are \(0\) and empty products are \(1\).  Bounded sums and products are
primitive recursive by ordinary recursion on \(b\).
\end{construction}

\begin{theorem}[Bounded Booleanization]
\label{thm:booleanization}
For every bounded formula \(A(\vec x)\), \(\iPRA\) derives uniformly:
\begin{align}
  &\chi_A(\vec x)=0\ \vee\ \chi_A(\vec x)=1,\label{eq:boolean-range}\\
  &A(\vec x)\ \longrightarrow\ \chi_A(\vec x)=1,\label{eq:to-char}\\
  &\chi_A(\vec x)=1\ \longrightarrow\ A(\vec x).\label{eq:from-char}
\end{align}
Consequently \(A\) is interchangeable with the quantifier-free equation
\(\chi_A=1\) in any finite trace invariant.
\end{theorem}

\begin{proof}
Proceed by structural induction on \(A\).  Equality and falsity are immediate from the
equality test and the successor axiom \(0\ne1\).  For conjunction, the induction
hypotheses make both factors Boolean.  Their product is \(1\) exactly when both are
\(1\); if either is \(0\), the product is \(0\).  For disjunction, the sign of the sum
is \(1\) exactly when at least one Boolean summand is \(1\).  These arguments use only
finite case analysis on numerical Booleans.

For implication, split on the Boolean value of \(\chi_A\).  If it is \(0\), the
characteristic value is \(1\), and \(A\to B\) follows because \(A\) would imply
\(\chi_A=1\), contradicting \(0=1\).  If \(\chi_A=1\), the characteristic value equals
\(\chi_B\), so the desired equivalence follows from the induction hypotheses for
\(B\).  This is an intuitionistically valid proof because the split is on a decidable
numerical equality, not on the arbitrary formula \(A\).

For the existential case, use induction on the bound \(b\) to prove that the bounded
sum is zero exactly when all characteristic values below \(b\) are zero, and has
positive sign exactly when some value is one.  The forward direction chooses the least
index whose summand is nonzero by bounded search; because each summand is Boolean, it
is one.  The reverse direction is immediate because a one summand makes the sum
positive.  Apply the induction hypothesis for the body to move between
\(\chi_A(u)=1\) and \(A(u)\).

For the universal case, induction on \(b\) proves that the bounded product is one
exactly when every factor below \(b\) is one.  The successor equation is
\[
   \prod_{u<b+1}\chi_A(u)
   =\left(\prod_{u<b}\chi_A(u)\right)\chi_A(b).
\]
A product of two Booleans is one exactly when both factors are one.  Again apply the
body induction hypothesis.  All inductions just mentioned have quantifier-free
invariants after replacing their bounded assertions by the corresponding sum or
product equations.  Thus they are available in \(\iPRA\).
\end{proof}

\subsection{Substitution and compilation discipline}

For use in a proof compiler, \cref{thm:booleanization} must be stable under
substitution.  If \(\sigma\) maps free variables to object terms, structural
construction gives
\[
   \chi_{A[\sigma]}=\chi_A[\sigma]
\]
up to the defining equations of the generated primitive-recursive symbols.  The Lean
implementation should prove this simultaneously with Booleanization, because the
bounded quantifier cases shift de Bruijn indices and are otherwise a common source of
errors.

The proposed \code{BForm} datatype in
\code{lean/FRNF/ObjectTargets.lean} records the grammar, with fields
\code{toFm} and \code{char}.  The three intended endpoints are exactly
\cref{eq:boolean-range,eq:to-char,eq:from-char}.  Completing them should reuse the
checked bounded-sum infrastructure rather than construct a second arithmetic library.

\subsection{Why implication is included}

One could normalize every bounded formula to positive Boolean combinations of atomic
equalities, but keeping implication is useful.  Contracts frequently have the form
``if the index is in range, then the table entry is valid.''  Compiling implication
directly avoids a separate negation-normal-form pass.  The construction remains
constructive because the antecedent has already been compiled to a Boolean.

\subsection{Consequence for induction}

Let \(I(t,\vec x)\) be a bounded assertion about the state of a loop at time \(t\).
Replace it by
\[
   \chi_I(t,\vec x)=1.
\]
The latter is a literal equation between primitive-recursive terms and hence a
quantifier-free formula.  If object derivations establish the base and successor
steps, the \(\PRAAx\) quantifier-free induction axiom closes the loop.  This is the
basic bridge from finite mathematics to the restricted induction of \(\iPRA\).

\section{Least bounded search and finite choice}
\label{sec:search}

Bounded search is the canonical way to turn a decidable finite existence statement
into a primitive-recursive witness.  It is also the basic deterministic selector used
inside demand passes, viability tables, finite nets, and certificate checkers.  The
repository already contains a checked object implementation in
\code{IPRADerivations.Core.Search}; we state the abstract theorem because it is a
building block of the normal form.

\subsection{Definition by primitive recursion}

Let \(p(i,x)\in\B\) be a primitive-recursive Boolean predicate, with the convention
that \(p(i,x)=1\) means success.  Define
\[
 \mu_p(0,x)=0
\]
and, for \(N\ge0\),
\[
 \mu_p(N+1,x)=
 \begin{cases}
   \mu_p(N,x),&\mu_p(N,x)<N,\\
   N,&\mu_p(N,x)=N\ \text{and}\ p(N,x)=1,\\
   N+1,&\mu_p(N,x)=N\ \text{and}\ p(N,x)=0.
 \end{cases}
\]
Every branch is expressible by equality tests and selectors, so \(\mu_p\) is primitive
recursive.  Equivalently, using a sign convention in which zero means success, one may
use the recurrence already implemented in the repository.

\begin{theorem}[Least bounded search]
\label{thm:bounded-search}
The function \(\mu_p\) satisfies, and \(\iPRA\) proves:
\begin{enumerate}[label=(\roman*),leftmargin=2.4em]
\item \(\mu_p(N,x)\le N\);
\item if \(\mu_p(N,x)<N\), then \(p(\mu_p(N,x),x)=1\);
\item if \(i<\mu_p(N,x)\), then \(p(i,x)=0\);
\item if \(i<N\) and \(p(i,x)=1\), then
      \(\mu_p(N,x)<N\) and \(\mu_p(N,x)\le i\).
\end{enumerate}
\end{theorem}

\begin{proof}
Use one induction on \(N\) with a joint quantifier-free invariant encoding (i)--(iii)
for an arbitrary comparison index \(i\).  The base case is immediate.  For the step,
there are two cases.  If the old search value is below \(N\), the recurrence freezes
it; all three clauses follow from the induction hypothesis.  If the old value equals
\(N\), inspect the Boolean \(p(N,x)\).  A success freezes the new value at \(N\); a
failure moves it to \(N+1\).  In either case bounded minimality follows by combining
the old minimality with the value at \(N\).  Clause (iv) is then derived by
contradiction: if \(i<\mu_p(N,x)\), clause (iii) contradicts \(p(i,x)=1\); otherwise
\(\mu_p(N,x)\le i<N\), and clause (ii) supplies success.  All case distinctions are on
numerical equalities or Boolean values, so the invariant is quantifier-free.
\end{proof}

\subsection{Finite choice is harmless when its bound is explicit}

A common constructive concern is the use of ``choose an \(i\)'' from a finite set.
In this setting the concern is not finite choice itself.  Once a bound \(N\) and a
Boolean verifier are available, \cref{thm:bounded-search} gives a canonical least
witness.  The real questions are:

\begin{itemize}
\item where the bound \(N\) came from;
\item whether the verifier captures the intended mathematical property; and
\item whether choosing the least encoded witness is extensionally legitimate.
\end{itemize}

The first question leads to producer boundaries and the cardinality-to-cutoff barrier.
The second is an object-proof obligation.  The third motivates relational outputs:
least choice is computationally available but may not respect equality of represented
objects.  One should therefore choose least only after proving that the downstream
observable is independent of the enumeration, or use least choice merely as an
internal path selector whose output is later erased by a stability theorem.

\subsection{Bounded searches over coded tuples}

Pairing and simultaneous recursion allow a search over any fixed finite product to be
reduced to one numerical interval.  For example, to search
\((i,j)\) with \(i<I\) and \(j<J\), use the code \(k=iJ+j<IJ\), recover
\(i=k\mathbin{/}J\) and \(j=k\bmod J\), and Booleanize the range conditions.  More
complex bounded records may be encoded by nested pairing.  The size of the search
space may be enormous; \(\iPRA\) is concerned with primitive recursiveness and proof,
not practical efficiency.  Nevertheless, the normal form encourages structured
loops and dynamic programming when they yield smaller invariants.

\subsection{The selector/finality distinction}

Suppose a finite cloud \(C_d(x)\) contains several valid approximations at precision
\(d\).  Bounded search yields the least valid code \(c_d\).  This proves existence of a
primitive-recursive selector for each finite level.  It does \emph{not} prove that
\((c_d)_d\) is coherent, convergent, or presentation invariant.  Those are separate
statements about compatibility edges, diameter bounds, and motion budgets.  We will
therefore use bounded search in two roles:

\begin{enumerate}[label=(\alph*),leftmargin=2.4em]
\item as an exact selector when equality and ordering of candidates have mathematical
meaning; and
\item as a deterministic implementation detail inside backward viability, followed by
a proof that the final observable does not depend on that detail.
\end{enumerate}

This distinction prevents a finite-choice theorem from being silently upgraded to a
choice principle for infinite represented objects.

\section{Static specialization and the universal-evaluator barrier}
\label{sec:static}

Many informal frameworks say that a mathematical object is supplied by ``a code for a
primitive-recursive function'' and then evaluate that code inside the theory.  This is
unsafe unless the code language is restricted or the evaluator is externally fixed.
The obstacle is elementary and exact.

\subsection{No total primitive-recursive universal evaluator}

\begin{barrier}[Diagonal obstruction]
\label{barrier:no-universal-pr}
There is no binary primitive-recursive function \(U(e,x)\) such that for every unary
primitive-recursive function \(f\) there is a code \(e\) with
\(U(e,x)=f(x)\) for all \(x\).
\end{barrier}

\begin{proof}
Assume such a \(U\) exists.  Then
\[
    g(x)=U(x,x)+1
\]
is primitive recursive, since the class of primitive-recursive functions is closed
under composition and successor.  By universality, choose \(e\) such that
\(U(e,x)=g(x)\) for all \(x\).  Evaluating at \(e\) gives
\[
    U(e,e)=g(e)=U(e,e)+1,
\]
a contradiction in \(\N\).  Hence no such \(U\) exists.
\end{proof}

The result does not say that primitive-recursive syntax cannot be coded, or that no
interpreter can evaluate bounded fragments.  It says that a total primitive-recursive
interpreter cannot be extensionally universal for all primitive-recursive programs.
A standard evaluator for primitive-recursive codes is partial recursive, and a
fuel-bounded evaluator is primitive recursive but may return ``unfinished.''

\subsection{Three admissible representation patterns}

\begin{definition}[Static specialization]
A fixed primitive-recursive evaluator \(f\) is represented in the object theory by its
own function symbol, defining equations, and derived contract.  A metatheoretic
compiler may map the external description of \(f\) to this syntax, but the object
program does not receive the description as data.
\end{definition}

This yields three sound patterns.

\paragraph{Fixed evaluator.}
A theorem is uniform in numerical inputs and precision but uses one externally fixed
coefficient-name evaluator, oracle, or transition function.  The compiler generates a
symbol for that evaluator.  This is the strongest and simplest pattern.

\paragraph{Finite table.}
For a fixed final demand, the reverse horizon pass determines a finite set of required
queries.  Their answers are supplied as a coded table satisfying a finite validity
formula.  The theorem never evaluates a program code at all.  This pattern is
particularly useful for extensional composition: only the rows actually consumed by
the proof are exposed.

\paragraph{Restricted interpreter.}
A fixed syntax class with an explicit structural or numerical fuel is interpreted by a
primitive-recursive evaluator.  The contract includes a proof that the supplied code
belongs to the class and that the fuel suffices.  Examples include straight-line
arithmetic expressions of bounded depth, fixed-order recurrences, or a bounded
instruction machine.

\subsection{Static families versus internal quantification}

Suppose \(F_e\) is a family indexed externally by primitive-recursive source code.  A
valid theorem schema may say:

\begin{quote}
For every externally fixed source \(e\), the compiler produces a symbol
\(\widehat F_e\), an object formula \(C_e\), and a derivation
\(\vdash_{\PRAAx} C_e\).
\end{quote}

This is a computable metatheoretic transformation if the compiler itself terminates.
It is not the same as one object formula quantifying over \(e\) and invoking a universal
\(U(e,x)\).  The distinction is analogous to the difference between an axiom schema
and a single truth predicate.

\subsection{Application: variable-degree polynomial constructions}

A polynomial degree \(n\) is ordinary numerical data, so a single primitive-recursive
algorithm may recurse on \(n\).  The subtlety arises when each coefficient is itself
specified by an arbitrary primitive-recursive evaluator passed by code.  The admissible
uniform statement is one of the following:

\begin{enumerate}[label=(\roman*),leftmargin=2.4em]
\item for each fixed coefficient evaluator, one factorization program works for every
degree and requested precision;
\item a finite coefficient table through the computed justification horizon suffices;
or
\item a restricted coefficient language is interpreted with explicit fuel.
\end{enumerate}

The first option closes the practical variable-degree composition gap while respecting
\cref{barrier:no-universal-pr}.  It is stronger than treating every degree as a
separate theorem but weaker---necessarily so---than an internal total evaluator for
all primitive-recursive codes.

\subsection{Lean architecture}

Static specialization should happen before the object syntax is constructed.  A Lean
metaprogram or ordinary definition receives a host representation of a component and
returns:

\begin{enumerate}[label=(\arabic*),leftmargin=2.4em]
\item a concrete \code{PR n} term;
\item its defining equations as applications of the existing \(\PRAAx\) schemes;
\item the object contract formula; and
\item a proof constructor specialized to that syntax.
\end{enumerate}

The host compiler may be generic and higher-order.  The generated object theorem is
first-order and static.  This is precisely where higher-type convenience is harmless:
it constructs syntax but does not become a numerical oracle inside \(\iPRA\).

\section{Dual horizons: computation versus justification}
\label{sec:dual-horizons}

\subsection{The under-specification caused by one horizon}

Let \(A\) be a represented object whose finite rows are \(A(0),A(1),\ldots\).  A
transducer \(F\) produces output row \(F(A)(d)\).  It is customary to give a use
function \(u_F(d)\) such that the program reads only rows at indices at most
\(u_F(d)\).  This is enough to run the program, but not necessarily enough to prove
its result valid.

For example, an output approximation may be computed from \(A(d+2)\), while its error
proof compares \(A(d+2)\) with an anchor \(A(2d+7)\).  Or a root routine may read a
rounded coefficient row but justify deflation by checking residuals at a shifted,
higher precision.  The proof horizon is then larger than the computation horizon.

\begin{definition}[Dual horizon]
A finite-use interface carries two monotone primitive-recursive functions
\[
    r_F,j_F:\N\to\N
\]
with \(r_F(d)\le j_F(d)\).  The computation of output row \(d\) reads input rows only
through \(r_F(d)\).  The embedded proof of the row contract uses input contract
instances only through \(j_F(d)\).
\end{definition}

Monotonicity can always be enforced by prefix closure:
\[
    \overline h(d)=\max_{q\le d}h(q).
\]
We therefore do not require every handwritten local horizon to be monotone; the
compiler normalizes it.

\subsection{Sequential composition}

Suppose \(F:X\to Y\) and \(G:Y\to Z\) have dual horizons.  To compute row \(d\) of
\(G(F(x))\), the outer routine may ask for any row \(q\le r_G(d)\) of \(F(x)\).  The
input rows of \(x\) needed for all such queries are bounded by
\[
   r_{G\circ F}(d)
   =\max_{q\le r_G(d)}r_F(q).
\]
For justification, the proof of the outer contract may instantiate input contracts
for any \(q\le j_G(d)\).  Each such inner contract requires rows through \(j_F(q)\), so
\[
   j_{G\circ F}(d)
   =\max_{q\le j_G(d)}j_F(q).
\]
If horizons are prefix-monotone, these simplify to
\(r_F(r_G(d))\) and \(j_F(j_G(d))\), but the prefix-maximum form is safer and is what
the reference implementation uses.

\begin{theorem}[Dual-horizon composition]
\label{thm:dual-compose}
If \(F\) and \(G\) satisfy their finite-use and correctness contracts with dual
horizons \((r_F,j_F)\) and \((r_G,j_G)\), then \(G\circ F\) satisfies them with the
horizons above.  Moreover
\[
   r_{G\circ F}(d)\le j_{G\circ F}(d).
\]
\end{theorem}

\begin{proof}
Every query issued by the computation of \(G\) at demand \(d\) has index
\(q\le r_G(d)\).  By the finite-use contract for \(F\), computing that row reads no
input above \(r_F(q)\).  Taking the finite maximum over all possible \(q\) proves the
computation bound.

For correctness, the local proof for \(G\) uses only validity instances of rows
\(q\le j_G(d)\).  The local proof for each such row of \(F\) uses input validity
instances only through \(j_F(q)\).  Simultaneous cut substitutes these inner
proofs into the outer proof, so no input instance above the displayed maximum occurs.

Finally, \(r_G(d)\le j_G(d)\).  For each \(q\le r_G(d)\), we have
\(r_F(q)\le j_F(q)\), and the latter value is included in the maximum defining
\(j_{G\circ F}(d)\).  Hence the computation maximum is at most the justification
maximum.
\end{proof}

\subsection{Multiple inputs and asymmetric demands}

A component with inputs \(X_1,\ldots,X_k\) carries one dual horizon per input.  Some
inputs may be exact parameters with constant horizon zero; others may be names queried
at very different rates.  Composition is performed componentwise.  This matters in
analysis, where a function name, a derivative modulus, a contour, and an error budget
have distinct demand behavior.

For a graph edge \(v\to u\), write
\[
   r_{v,u}(d),\qquad j_{v,u}(d)
\]
for the demand placed on dependency \(u\) when node \(v\) is required through row
\(d\).  The graph-level closure in \cref{sec:demand-graphs} joins demands arriving
from several users by componentwise maximum.

\subsection{Proof horizons as a noninterference theorem}

The justification horizon has a useful extensional reading.  If two input
presentations agree on rows and validity certificates through \(j_F(d)\), then the
compiled proof of output row \(d\) is the same up to substitution of equal data.  Thus
\(j_F(d)\) is not just an implementation bound; it is a finite-observation theorem
about the proof itself.

This permits modular replacement of infinite names by finite tables.  To establish an
output contract at demand \(d\), it suffices to supply a finite table through
\(j_F(d)\) satisfying the finitely many input contract instances.  No global object
representing the entire infinite name need occur inside the induction.

\subsection{Adaptive queries}

If the next query index depends on previous answers, a simple syntactic scan does not
produce a fixed list.  The solution is to bound the adaptive computation by a clock and
compute a worst-case demand transformer.  Let a state contain the current query index,
answer table, and remaining fuel.  A primitive-recursive majorant \(H(d,t)\) bounds
all queries reachable in \(t\) steps from final demand \(d\).  The dual horizon is the
prefix maximum of this majorant over the clock, with a separate recurrence for proof
queries.  The finite trace then records the actual path, while the precomputed horizon
makes the induction independent of an infinite oracle.

\subsection{Reference artifact}

The class \code{DualHorizon} in \code{reference/frnf.py} implements validation and
sequential composition.  The regression tests include a case where the justification
horizon is strictly larger than the read horizon.  The host-level Lean record in
\code{lean/FRNF/Semantic.lean} mirrors the definition.  Neither artifact by itself is
an object proof; the object implementation is obtained after horizons are represented
by fixed \code{PR} symbols and their finite maxima are derived in \(\Prf\,\PRAAx\).

\section{Finite demand graphs and reverse closure}
\label{sec:demand-graphs}

Real proofs are rarely simple chains.  A requested observable may depend on several
approximations, each produced from shared lower-level data, with proof obligations
that fan out differently from computation.  We therefore organize a proof as a fixed
acyclic graph.

\subsection{Graph data}

Let \(V=\{0,\ldots,m-1\}\) be nodes in topological order: every dependency of node
\(v\) has smaller index.  For each edge \(v\to u\), where \(u\) is a dependency of
\(v\), supply monotone dual demand transformers
\[
  r_{v,u},j_{v,u}:\N\to\N,
  \qquad r_{v,u}(d)\le j_{v,u}(d).
\]
A demanded node carries a pair \((R_v,J_v)\).  We adjoin a bottom value
\(\bot\), meaning ``this node is not used.''  This distinction is essential: numerical
demand \(0\) means that row zero is required, not that the node is unused.  Bottom is
below every pair; nonbottom pairs are ordered componentwise and joined by maximum.

Given a root \(z\) and requested output precision \(d\), initialize
\((R_z,J_z)=(d,d)\) and every other node to \(\bot\).  Traverse nodes in reverse
topological order.  Skip bottom nodes.  For each edge \(v\to u\) from a nonbottom node,
update
\[
\begin{split}
 R_u&\leftarrow\max\left(R_u,\max_{q\le R_v}r_{v,u}(q)\right),\\
 J_u&\leftarrow\max\left(J_u,\max_{q\le J_v}j_{v,u}(q)\right).
\end{split}
\]
Here joining a propagated pair into \(\bot\) returns that pair.  For multiple roots,
initialize each requested root and take joins.

\begin{theorem}[Reverse demand closure]
\label{thm:demand-closure}
The reverse pass computes the least bottom-extended demand assignment closed under
all edges leaving demanded nodes and dominating the root demand.  Its nonbottom nodes
are exactly the dependency cone of the root.  It is primitive recursive in the fixed
graph data and the final demand.  Every resulting pair satisfies \(R_v\le J_v\).
\end{theorem}

\begin{proof}
Because the order is topological, when node \(v\) is processed in reverse order, all
users of \(v\) have already propagated their final demands to it.  No later step can
increase \(v\)'s demand.  A node receives a nonbottom demand exactly when it lies on a
dependency path from the root.  Therefore one pass suffices and disconnected nodes do
not spuriously request row zero.

Closure follows directly from the update equations.  To prove minimality, let
\((R'_v,J'_v)\) be any closed assignment dominating the root.  Descend through the
reverse order.  Assume every already processed node has computed demand at most its
primed demand.  Monotonicity of prefix maxima and closure of the primed assignment show
that every propagation from the current node is at most the corresponding primed
dependency demand.  Taking maxima preserves the inequality.  Induction proves
\((R_v,J_v)\le(R'_v,J'_v)\) for all \(v\).

Primitive recursiveness follows because the graph is fixed and every update is a
bounded maximum of primitive-recursive functions.  If the graph itself is numerical
input, encode its finite adjacency table and run a bounded loop of length \(m\); the
same argument becomes a quantifier-free invariant about the table.

For \(R_v\le J_v\), induct on the reverse pass.  Initially the property holds.  On an
edge, \(R_v\le J_v\), the edge read transformer is bounded by its justification
transformer, and prefix maximum is monotone.  Hence the propagated read demand is at
most the propagated justification demand; componentwise joins preserve the property.
\end{proof}

\subsection{Why closure precedes induction}

The demand pass is run in the metatheory or as a preliminary primitive-recursive
computation, before the main correctness induction.  Its output is a finite vector of
bounds.  The proof then materializes exactly the table rows and contract instances
within those bounds.  This order avoids an illegal invariant saying ``for every future
precision, all needed rows are valid.''

In an object implementation, the demand pass itself is a bounded loop and may be
verified by quantifier-free induction.  The main trace uses its completed table as
fixed input.  Nested uses of quantifier-free induction are harmless: each invariant is
finite and explicit.

\subsection{Shared dependencies and proof reuse}

If two nodes depend on the same source, the closure takes the maximum of their demands.
The source table is computed once.  Likewise, its validity derivations may be stored
once in the proof schedule and cut into both consumers.  This is the proof-theoretic
analogue of common-subexpression elimination.

A graph representation is therefore not merely convenient notation.  It can reduce
proof terms and makes dependency auditing possible.  Each node declares exactly which
earlier formulas its local proof uses; undeclared global assumptions cannot enter
silently.

\subsection{Cycles and feedback}

A genuine cyclic dependency cannot be handled by a topological pass.  In practice one
of three situations occurs.

\begin{enumerate}[label=(\roman*),leftmargin=2.4em]
\item The apparent cycle is a bounded iteration.  Collapse it into one clocked loop
node whose state contains all mutually recursive values.
\item The cycle is a fixed-point specification with a contraction modulus.  Replace it
by a finite Picard iteration whose clock is derived from the modulus.
\item The cycle expresses an infinitary or impredicative construction with no supplied
majorant.  It lies outside the unconditional \(\FRNF\) interface and must remain a
boundary.
\end{enumerate}

Simultaneous primitive recursion, already available in the checked repository library,
is the standard implementation of the first case.

\subsection{Object encoding}

For a fixed small graph, the cleanest formalization generates one tuple of demand terms
per node and proves the closure equations directly.  For a graph supplied as finite
data, encode:

\begin{itemize}
\item the node count and topological order;
\item edge offsets in a flattened adjacency table;
\item codes selecting from a fixed finite family of demand transformers;
\item an active bit distinguishing bottom from the genuine numerical demand zero; and
\item the read/justification arrays updated by primitive recursion.
\end{itemize}

The interpreter is admissible because the family of node and edge operations is fixed
and finite; it is not a universal evaluator for arbitrary primitive-recursive syntax.
The object theorem \code{demandClosure_correct} in the Lean roadmap is the intended
closed derivation of this finite table algorithm.

\section{Clocked traces and quantifier-free internalization}
\label{sec:trace}

Once all finite demands are known, every adaptive construction is represented as a
clocked transition system.  This section gives the central object-proof pattern.

\subsection{Trace specification}

Fix primitive-recursive symbols:
\[
  \operatorname{init}(x),\quad
  \operatorname{step}(t,s,x),\quad
  \operatorname{fuel}(x),\quad
  \operatorname{out}(s,x).
\]
Define the state sequence by primitive recursion:
\[
  S(0,x)=\operatorname{init}(x),\qquad
  S(t+1,x)=\operatorname{step}(t,S(t,x),x).
\]
The state is a numerical code for every accumulator, table, flag, and local certificate
needed by the process.  Let \(I(t,s,x)\) be a bounded formula expressing the invariant,
and let \(Q(x,\operatorname{out}(S(F,x),x))\) be the desired terminal contract, where
\(F=\operatorname{fuel}(x)\).

Booleanize \(I\) and abbreviate
\[
  \widehat I(t,x)\;:\equiv\;\chi_I(t,S(t,x),x)=1.
\]
This is a quantifier-free equation.

\begin{theorem}[QF trace internalization]
\label{thm:trace-internalization}
Suppose \(\iPRA\) derives:
\begin{align*}
 &\widehat I(0,x),\\
 &\widehat I(t,x)\longrightarrow\widehat I(t+1,x),\\
 &\widehat I(F,x)\longrightarrow
   Q(x,\operatorname{out}(S(F,x),x)).
\end{align*}
Then \(\iPRA\) derives
\[
  Q(x,\operatorname{out}(S(F,x),x)).
\]
The construction is uniform in the fixed primitive-recursive syntax and yields an
embedded \(\Prf\,\PRAAx\) term from embedded base, step, and finish derivations.
\end{theorem}

\begin{proof}
Apply quantifier-free induction to \(\widehat I(t,x)\), using the first two derivations.
Universal introduction closes the induction variable, and universal elimination
specializes the result at the term \(F\).  Modus ponens with the terminal derivation
yields the output contract.  At the level of the embedded calculus, these are exactly
\code{PRA.qfInd}, \code{Prf.allE}, \code{Prf.impI}, and \code{Prf.impE}.  No semantic
soundness theorem is needed to construct the proof term.
\end{proof}

The structure \code{QFLoop} and theorems \code{close}, \code{at}, and \code{finish} in
\code{lean/FRNF/Kernel.lean} are a direct source-level implementation of this proof.

\subsection{What belongs in the state?}

A robust state code contains enough information that the step invariant is local.  A
typical record includes:

\begin{itemize}
\item current clock and active/terminal flag;
\item current numerical approximation or exact accumulator;
\item a finite input table or indices into a shared table;
\item a running maximum of requested proof rows;
\item a rank or remaining fuel;
\item local residual, error numerator, or divisibility certificate;
\item finite candidate masks for relational stages; and
\item a checksum or defining-equation witness when useful for efficient reflection.
\end{itemize}

Including more fields can make the invariant simpler.  Since all records are encoded by
pairing and projections, this is a proof-engineering tradeoff rather than a change of
logical strength.

\subsection{Frozen terminal states}

Every clocked process should be made total by freezing after termination:
\[
  \operatorname{step}(t,s,x)=s
  \quad\text{when}\quad\operatorname{terminal}(s,x)=1.
\]
This eliminates partiality and allows a fixed clock larger than the actual stopping
time.  The invariant includes ``terminal states remain terminal and preserve their
output.''  Majorized descent in \cref{sec:descent} supplies a common source of such
clocks.

\subsection{Local reflection versus global semantics}

The trace theorem requires embedded proofs of the base and step equations.  A host
proof that the transition is correct can guide these derivations but cannot replace
them.  In practice the object step proof is assembled from:

\begin{enumerate}[label=(\roman*),leftmargin=2.4em]
\item defining equations for the primitive-recursive transition;
\item Booleanization lemmas for branch guards;
\item arithmetic library lemmas for the updated fields; and
\item local component contracts already present in the proof schedule.
\end{enumerate}

This decomposition is favorable to automation.  The transition itself is fixed and
finite; the compiler can generate the formula and most congruence steps.  Domain
lemmas remain explicit.

\subsection{Variable-length recursion}

A loop whose length is a numerical input is not a schema of separate finite proofs.
It is one primitive recursion and one induction on the clock.  The final theorem may
therefore be genuinely uniform in a matrix size, polynomial degree, path length, or
search bound.  The restriction is that the transition and invariant must be fixed
primitive-recursive syntax.  This is how variable-degree constructions avoid the
universal-evaluator barrier.

\subsection{Nested finite loops}

Nested loops can be flattened into one state machine by storing both counters, or
formalized modularly as components whose local contracts are composed.  Flattening
uses one induction but may produce a large invariant; modularization creates smaller
proofs but requires a proof schedule.  \(\FRNF\) supports both.  The preferred choice
is to isolate mathematically meaningful subroutines and flatten only tightly coupled
adaptive loops.

\section{Proof-producing schedules and simultaneous cut}
\label{sec:proof-schedule}

The trace compiler turns one bounded process into an object derivation.  A complete
mathematical proof contains many such components.  The proof schedule is the final
logical assembly layer.

\subsection{Schedule invariant}

Fix a permanent context \(\Gamma\) of input and boundary assumptions.  Let
\(D=[C_k,\ldots,C_1]\) be a reverse-topological list of contracts already compiled.
A schedule carries the invariant
\[
  \forall C\in D,\qquad \Gamma\vdash_{\PRAAx} C.
\]
To add a new node \(A\), provide a dependency sublist \(\Delta\subseteq D\) and a local
derivation
\[
  \Delta,\Gamma\vdash_{\PRAAx}A.
\]
Simultaneous cut replaces each dependency hypothesis by the corresponding stored
derivation, and each permanent hypothesis by itself.  The result is
\(\Gamma\vdash_{\PRAAx}A\).

\begin{theorem}[Finite proof schedule]
\label{thm:proof-schedule}
A finite topologically ordered list of local embedded derivations compiles to embedded
derivations of all node contracts from the permanent context.  In particular, the
root contract is derivable from exactly the listed boundary assumptions.
\end{theorem}

\begin{proof}
Induct on the node list.  The empty schedule is vacuous.  At a successor, apply the
simultaneous substitution/cut operation of the embedded deduction system to the local
proof.  The induction hypothesis supplies derivations for declared dependencies, and
ordinary hypotheses supply the permanent context.  Preserve all earlier derivations
unchanged.  The root is retrieved by membership.
\end{proof}

The \code{Schedule.extend} definition in the accompanying Lean source implements this
argument using \code{Prf.trans}.  It contains no semantic cast and no appeal to
soundness.

\subsection{Dependency declarations as an audit mechanism}

A local proof is accepted only under \(\Delta++\Gamma\), not under the entire global
collection of facts.  The inclusion proof \(\Delta\subseteq D\) is the acyclicity
certificate.  This has three benefits.

\begin{enumerate}[label=(\roman*),leftmargin=2.4em]
\item The generated proof term records the actual dependency graph.
\item A boundary assumption cannot leak into a supposedly unconditional component
unless it appears in \(\Gamma\) and the theorem statement.
\item Components can be replaced or strengthened locally, and unused earlier nodes do
not enter the final proof.
\end{enumerate}

For a generated library, one can additionally emit a machine-readable manifest mapping
each node to source files, theorem names, horizon transformers, and status.

\subsection{Proof-producing compilation versus proof checking}

A conventional certificate checker takes a numerical certificate and returns true or
false.  A proof-producing compiler instead constructs a syntax tree of deduction
rules.  The two may share the same finite algorithm, but their trust boundaries differ.

In the present architecture, the generated proof term can be large.  This is not a
logical defect.  Compression can be addressed later by factoring reusable lemmas,
using reflection theorems already proved in the object theory, or exporting to the
smaller native scalar calculus.  The first goal is a transparent proof with the right
axioms.

\subsection{Closed theorems and conditional theorems}

If \(\Gamma=[]\), the schedule yields a closed theorem.  If \(\Gamma\) contains a
modulus, gap, height bound, or other producer contract, the schedule yields a precise
conditional theorem.  When a producer is later proved, the schedules compose by one
more cut.  This makes partial progress cumulative: formalizing the finite verifier is
valuable even before the deep producer theorem is available.

\subsection{Native export}

After the ordinary derivation is complete, the repository theorem
\code{pra_iff_native} translates it to the native scalar core.  The final export stage
is therefore
\[
  \Gamma\vdash_{\PRAAx}\varphi
  \quad\Longrightarrow\quad
  \operatorname{ofPRA}(\Gamma)\vdash_{\mathrm{native}}
  \operatorname{ofPRA}(\varphi).
\]
This translation should occur only at the end.  The ordinary presentation has a much
richer checked arithmetic library, while the native core is the compact foundational
endpoint.

\subsection{A generated node template}

A practical compiler may emit one Lean declaration per node in the following pattern.

\begin{lstlisting}
-- Formula and fixed PR symbols for node v.
def nodeVFormula : Fm PTm n := ...

-- Local proof: only declared predecessors and base assumptions.
theorem nodeV_local : depsV ++ Gamma |-[(PRAAx)] nodeVFormula := by
  -- unfold fixed transition equations
  -- invoke Booleanization / bounded-search / arithmetic lemmas
  -- close a QF loop if needed
  ...

-- Add to proof schedule by simultaneous cut.
def scheduleV : Schedule Gamma (nodeVFormula :: done) :=
  Schedule.extend schedulePrev deps_subset nodeV_local
\end{lstlisting}

The notation in actual source follows the repository's scoped syntax.  The key property
is the type of \code{nodeV_local}, not the tactic script used to construct it.

\section{Finite relational witnesses and backward viability}
\label{sec:viability}

A deterministic witness is often the wrong intermediate object.  At finite precision,
there may be several equally valid roots, charts, lifts, combinatorial fillings, or
approximate optimizers.  Forcing an early choice can make a correct construction
nonextensional or can smuggle in a discontinuous decision.  \(\FRNF\) therefore treats
a finite relation as a first-class output.

\subsection{Layered finite relations}

Fix a horizon \(h\).  At level \(i\le h\), candidates are numbered
\(a<B_i\).  Let \(V_i(a)\in\B\) be the validity predicate and let
\(E_i(a,b)\in\B\) express compatibility between candidate \(a\) at level \(i\)
and candidate \(b\) at level \(i+1\).  All bounds and predicates are primitive
recursive in the external parameters and final demand.

A complete path is a sequence \(a_0,\ldots,a_h\) with
\[
  V_i(a_i)=1\quad(i\le h),
  \qquad E_i(a_i,a_{i+1})=1\quad(i<h).
\]
The path need not be unique.

\subsection{Backward viability recurrence}

Define viability masks backward:
\[
\begin{split}
 W_h(a)&=V_h(a),\\
 W_i(a)&=V_i(a)\cdot
   \sg\!\left(\sum_{b<B_{i+1}}E_i(a,b)W_{i+1}(b)\right).
\end{split}
\]
Thus \(W_i(a)=1\) exactly when \(a\) can be extended to a valid path through level
\(h\).  After computing the masks, choose \(a_0\) as the least viable candidate and,
recursively, choose \(a_{i+1}\) as the least viable successor compatible with
\(a_i\).

\begin{theorem}[Backward viability]
\label{thm:backward-viability}
For every finite layered relation:
\begin{enumerate}[label=(\roman*),leftmargin=2.4em]
\item \(W_i(a)=1\) if and only if \(a\) extends to a valid path from level \(i\) to
level \(h\);
\item a complete path exists if and only if some \(a<B_0\) has \(W_0(a)=1\);
\item if a complete path exists, least bounded search applied to the masks and edges
constructs a complete path; and
\item the masks and selected path are primitive recursive and their correctness is
provable in \(\iPRA\) by quantifier-free induction.
\end{enumerate}
\end{theorem}

\begin{proof}
For (i), use backward induction on \(h-i\).  At level \(h\), viability is exactly
validity, which is the condition for the one-element suffix.  At level \(i<h\), the
recurrence says that \(a\) is valid and there exists a compatible \(b\) with
\(W_{i+1}(b)=1\).  By the induction hypothesis, that \(b\) extends to a valid suffix;
prepending \(a\) gives a valid path.  Conversely, the second element of any valid path
is a compatible viable successor, so the bounded sum has positive sign.

Clause (ii) is the case \(i=0\) of (i).  For (iii), choose the least viable initial
candidate.  If \(a_i\) is viable, the defining recurrence guarantees a compatible
viable successor; least bounded search finds one.  Induction constructs the complete
path and preserves validity and compatibility.

For (iv), encode the level-candidate table as a finite array.  The backward pass has
clock \(h+1\), and each mask entry is a bounded sum of Boolean products.  The forward
pass also has clock \(h+1\) and uses the checked least-search theorem.  The invariant
states, in Booleanized form, that processed masks satisfy the recurrence and the
selected prefix is valid, compatible, and viable.  Every quantifier is bounded by the
encoded table sizes and is therefore replaced by a characteristic equation using
\cref{thm:booleanization}.  Quantifier-free induction closes both passes.
\end{proof}

\subsection{Why no infinite choice is used}

The theorem is finite for each requested horizon.  It does not appeal to K\"onig's
lemma, dependent choice, or an infinite branch through the layered graph.  This is
crucial.  Many applications consume only a finite horizon, and stable-observer
arguments can show that the final result is independent of which finite path is used.
When an actual represented limit is required, coherence is proved quantitatively from
shrinking diameters and motion budgets, as in \cref{sec:stability}; it is not obtained
by selecting an abstract infinite branch.

\subsection{Relational composition}

Suppose a producer returns a finite relation \(R(x,a)\), while a consumer at state
\(s\) may accept several answers and move to \(s'=T(s,a)\).  A locally valid answer
can lead to a dead end at a later stage.  Backward viability solves this global finite
constraint problem: candidate \((s,a)\) is retained only when it has a compatible
continuation to the fixed final horizon.  The forward selector then produces a
successful interaction trace.

This is the correct finite replacement for a higher-type oracle strategy.  The entire
answer tree is finite after the final demand and clock are fixed.  The backward pass
computes a winning region; the forward pass selects one witness.  If the consumer's
observable is stable over all winning traces, the internal selection is erased.

\subsection{Applications}

The same recurrence covers several apparently different situations.

\begin{itemize}
\item \textbf{Root clouds.}  Candidates are rational disks or boxes with residual
certificates; edges express overlap or controlled movement.
\item \textbf{Charted geometry.}  Candidates are local charts or transition branches;
edges express overlap compatibility.
\item \textbf{Simplicial observation.}  Candidates are vertices in a carrier cloud;
edges express admissible selections.
\item \textbf{Approximate minimizers.}  Candidates are net points within a tolerance;
edges enforce motion between refinements.
\item \textbf{Finite gluing.}  Candidates are local sections; compatibility is equality
or witnessed equality on overlaps.
\item \textbf{Adaptive proof search.}  Candidates are local certificates; edges express
that the next verifier state accepts them.
\end{itemize}

The domain-specific work lies in defining validity and compatibility and proving
nonemptiness.  The viability compiler is generic.

\section{Stable observers, shrinking viability, and exact collapse}
\label{sec:stability}

Backward viability constructs one finite path, but the mathematical result should
usually not depend on the arbitrary least-code convention.  Stability converts a
relational witness into a canonical observable without claiming a canonical witness.

\subsection{Finite stable observers}

Let \(C\) be a nonempty finite candidate set and let \(O(c)\) be a rational or coded
observable.  Suppose
\[
   \forall c,c'\in C,
   \qquad d(O(c),O(c'))\le\eps.
\]
Then any bounded selector \(s(C)\in C\) yields an output within \(\eps\) of every other
valid output.  The selector may depend on enumeration, but its observable is stable.

\begin{theorem}[Stable-observer erasure]
\label{thm:stable-observer}
Assume embedded derivations certify that \(C_1,\ldots,C_m\) are valid finite
certificates and that their joint stability contract implies a goal \(G\).  Then
simultaneous cut yields an embedded derivation of \(G\) without exposing the internal
selectors.  In particular, if every valid candidate gives the same exact coded
observable, the result is exact.
\end{theorem}

\begin{proof}
The logical statement is immediate: from derivations of the certificate formulas and a
local derivation of \(G\) under those formulas, apply simultaneous cut.  The
mathematical burden is entirely in the stability derivation, which is domain-specific
but finite.  The \code{stableMany} theorem in the Lean kernel file implements the
logical step.
\end{proof}

\subsection{Finite-horizon coherence without an infinite branch}

For an actual represented output, we need a sequence of approximants.  Let the layered
relation extend indefinitely in the metatheory, but perform only finite viability
computations.  Choose a nondecreasing primitive-recursive horizon schedule
\(H(k)\ge k+1\).  Run backward viability through level \(H(k)\), select the least
viable path, and let \(x_k\) be the observable of its candidate at level \(k\).

Write \(V_{k,h}\) for the candidates at level \(k\) that extend to level \(h\).  These
sets shrink with \(h\).  Assume quantitative data:
\[
\begin{split}
 &\diam\{O_k(a):a\in V_{k,H(k)}\}\le\delta_k,\\
 &E_k(a,b)=1\quad\Longrightarrow\quad
   d(O_k(a),O_{k+1}(b))\le\eta_k.
\end{split}
\]

\begin{theorem}[Shrinking-viability coherence]
\label{thm:shrinking-viability}
Under the assumptions above,
\[
   d(x_k,x_{k+1})\le\delta_k+\eta_k.
\]
Consequently, if a primitive-recursive modulus \(M\) satisfies
\[
   \sum_{i=m}^{n-1}(\delta_i+\eta_i)\le 2^{-p}
   \quad\text{whenever}\quad n>m\ge M(p),
\]
then \((x_k)\) is a regular Cauchy name with modulus \(M\).  Its construction uses
only finite viability tables.
\end{theorem}

\begin{proof}
The path selected for horizon \(H(k+1)\) reaches at least level \(H(k)\), since the
schedule is nondecreasing.  Its candidate \(a'_k\) at level \(k\) therefore belongs to
\(V_{k,H(k)}\).  The candidate \(a_k\) used to define \(x_k\) also belongs to this set.
The diameter bound gives
\[
   d(O_k(a_k),O_k(a'_k))\le\delta_k.
\]
The candidate \(a'_{k+1}\) on the longer path is compatible with \(a'_k\), so the
motion bound gives
\[
   d(O_k(a'_k),O_{k+1}(a'_{k+1}))\le\eta_k.
\]
The latter observable is \(x_{k+1}\).  The triangle inequality proves the one-step
bound.  Summing from \(m\) to \(n-1\) and applying the supplied tail modulus proves the
Cauchy estimate.  Every approximant uses a finite table; no infinite branch is
selected.
\end{proof}

The theorem separates two sources of uncertainty.  \(\delta_k\) measures disagreement
among finite-horizon choices at the same level; \(\eta_k\) measures the movement
between adjacent levels along one compatible path.  Either can be zero in special
applications.

\subsection{Metastability is not a rate}

A finite-window theorem may say that for every \(p\) and control function \(g\), there
is \(N\le B(p,g)\) such that
\[
   d(x_i,x_j)\le2^{-p}
   \quad(N\le i,j\le N+g(N)).
\]
This is a useful finite contract, but it does not by itself give a tail modulus
\(M(p)\) valid for all \(i,j\ge M(p)\).  \(\FRNF\) therefore records metastability and
Cauchy rates as different interface types.  A downstream finite consumer may use the
window directly.  A represented limit requires an independent tail argument such as
the summable-motion bound in \cref{thm:shrinking-viability}.

\subsection{Exact branch collapse under separation}

Stable approximation does not automatically identify an exact discrete label.  A
separation certificate makes such a collapse legitimate.

\begin{theorem}[Gap collapse]
\label{thm:gap-collapse}
Let exact labels \(\lambda_0,\ldots,\lambda_{m-1}\) in a metric presentation satisfy
\[
   d(\lambda_i,\lambda_j)\ge\gamma>0
   \quad(i\ne j).
\]
If a candidate observation \(y\) is certified to lie within \(\gamma/3\) of some label,
then there is a unique such label.  Bounded search over the finite label list recovers
it, and all valid candidates recover the same label.
\end{theorem}

\begin{proof}
If \(y\) were within \(\gamma/3\) of two distinct labels, the triangle inequality
would give their distance at most \(2\gamma/3<\gamma\), contradicting separation.
Existence is assumed by the approximation certificate; uniqueness makes the least
bounded search independent of enumeration among valid witnesses.  The entire argument
is finite and can be Booleanized.
\end{proof}

For signs, charts, isolated roots, or simple eigenvalues, \(\gamma\) is an apartness or
spectral gap.  At collisions no such certificate exists, and the correct output is a
cluster, symmetric polynomial, spectral subspace, or another branch-invariant object.

\subsection{Equality versus observational equality}

Two finite witnesses need not have equal codes.  Stable-observer theorems derive the
equality or closeness of the observations actually consumed by later nodes.  This is a
form of proof-relevant extensionality: equality is established at the interface rather
than imposed by quotienting all presentations in the metatheory.  It is particularly
well suited to first-order arithmetic, where finite equality witnesses can be coded and
checked.

\section{Majorized descent and totalized termination}
\label{sec:descent}

Many classical proofs say ``repeat until no reduction is possible.''  To live in
\(\iPRA\), the process needs a primitive-recursive clock.  A decreasing numerical rank
provides one.

\subsection{Frozen descent}

Let \(S\) be a coded state space, \(T(s)\) a primitive-recursive terminal predicate,
\(N(s)\) a primitive-recursive next-state function, and
\(\rank(s)\in\N\) a primitive-recursive rank.  Define the totalized step
\[
  \widetilde N(s)=
  \begin{cases}
    s,&T(s)=1,\\
    N(s),&T(s)=0.
  \end{cases}
\]
Assume the local descent contract
\[
  T(s)=0\quad\Longrightarrow\quad \rank(N(s))<\rank(s).
\]

\begin{theorem}[Majorized descent]
\label{thm:majorized-descent}
Starting from \(s_0\), the frozen iteration
\(s_{t+1}=\widetilde N(s_t)\) is terminal by time \(\rank(s_0)+1\).  The terminal state
and full trace are primitive recursive, and the statement is provable in \(\iPRA\) by
quantifier-free induction.
\end{theorem}

\begin{proof}
Assume for contradiction that none of
\(s_0,\ldots,s_{\rank(s_0)+1}\) is terminal.  Then the rank strictly decreases at every
step, giving a descending sequence of \(\rank(s_0)+2\) natural numbers all at most
\(\rank(s_0)\).  More directly, prove by induction on \(t\) the bounded assertion
\[
  \bigl(\forall i<t\;T(s_i)=0\bigr)
  \longrightarrow
  \rank(s_t)+t\le\rank(s_0).
\]
After Booleanization this is quantifier-free.  At
\(t=\rank(s_0)+1\), nontermination would imply a natural number plus
\(\rank(s_0)+1\) is at most \(\rank(s_0)\), impossible.  Once a terminal state is
reached, freezing preserves it and its output.  Primitive recursiveness follows from
the fixed bounded iteration.
\end{proof}

\subsection{Lexicographic and weighted ranks}

A transition may decrease a tuple lexicographically rather than a single obvious
number.  If components have known bounds \(a_i<B_i\), encode the tuple in mixed radix:
\[
  \rho(a_0,\ldots,a_{k-1})
  =a_0\prod_{j>0}B_j+a_1\prod_{j>1}B_j+\cdots+a_{k-1}.
\]
A strict lexicographic decrease decreases \(\rho\).  If bounds vary with parameters,
they remain primitive-recursive terms.  Alternatively, use a weighted norm
\(\sum_i w_i a_i\) with weights chosen so that every decrease in an earlier component
dominates all possible increases in later ones.

The bound information is not optional.  An unbounded lexicographic tuple cannot in
general be compressed into a rank whose decrease is immediate.  The interface should
therefore store component bounds together with the transition.

\subsection{Euclidean algorithm}

For the state \((a,b)\) with \(b>0\), the Euclidean step is
\[
  (a,b)\longmapsto (b,a\bmod b).
\]
The second component strictly decreases, so \(\rank(a,b)=b\) is a rank.  Freezing at
\(b=0\) yields a clock of at most \(b_0+1\).  The checked ProveIt gcd construction uses
this pattern, and the reference Python implementation
\code{euclid_majorized} exposes its trace.

The extended Euclidean algorithm carries additional coefficient fields.  They may grow,
but termination is still controlled by the remainder rank.  The invariant relates the
coefficients to the current remainders; only the rank field must decrease.

\subsection{Descent proofs of impossibility}

Infinite descent arguments can be compiled in two related ways.

\begin{enumerate}[label=(\roman*),leftmargin=2.4em]
\item A witness transformation sends every nonminimal solution to a smaller solution.
Bounded search chooses the least solution below an explicit bound, producing a direct
contradiction.
\item A solution-count function below rank \(N\) is shown to vanish by induction on
\(N\): a solution at rank \(N+1\) would yield one at smaller rank, contradicting the
induction hypothesis.
\end{enumerate}

The repository's generic descent modules implement both styles.  \(\FRNF\) treats them
as reusable nodes rather than rebuilding descent for each theorem.

\subsection{When no rank is known}

A semantic proof that a process terminates does not automatically provide a
primitive-recursive clock provable in \(\iPRA\).  If the only evidence is an abstract
well-founded relation or a compactness argument, the process is not yet in \(\FRNF\).
One must supply a numerical majorant, restrict the domain, or expose termination as a
boundary assumption.  This is another instance of producer/verifier separation: a
step checker does not produce a stopping bound.

\section{Quantitative compactness through certified finite nets}
\label{sec:finite-nets}

Classical analysis often invokes compactness to pass from local information to a
uniform conclusion or to obtain an extremum.  In a finitary setting, compactness is an
interface consisting of a finite net and a modulus.  This section develops the generic
minimum theorem; the same pattern supports uniform continuity, covering arguments,
finite-dimensional topology, and integration.

\subsection{Net presentation}

Let \((X,d)\) be a represented metric object.  A certified net interface supplies, for
each precision parameter \(k\):

\begin{itemize}
\item a finite list \(N_k=(p_{k,0},\ldots,p_{k,m(k)-1})\);
\item a radius \(\delta(k)\in\Q_{>0}\); and
\item a finite certificate that every admissible input observation is within
\(\delta(k)\) of some net point.
\end{itemize}

For a function \(f:X\to\mathbb R\), a modulus supplies a rational error
\(\omega(k)\) such that
\[
  d(x,y)\le\delta(k)\quad\Longrightarrow\quad
  \abs{f(x)-f(y)}\le\omega(k).
\]
The evaluator returns rationals \(v_{k,i}\) with
\[
  \abs{v_{k,i}-f(p_{k,i})}\le\eta(k).
\]
All finite data and bounds are primitive recursive in \(k\).

\subsection{Finite minimum algorithm}

Choose the least index \(i_*\) minimizing \(v_{k,i}\).  This is a primitive-recursive
fold or bounded search.  Put \(m_k=v_{k,i_*}\) and \(p_k=p_{k,i_*}\).

\begin{theorem}[Finite-net minimum]
\label{thm:finite-net-min}
Under the net, modulus, and evaluator contracts:
\begin{enumerate}[label=(\roman*),leftmargin=2.4em]
\item for every admissible \(x\in X\),
\[
   f(p_k)\le f(x)+\omega(k)+2\eta(k);
\]
\item if \(\inf_X f\) is used in the external semantics, then
\[
   \inf_X f-\eta(k)\le m_k
   \le \inf_X f+\omega(k)+\eta(k);
\]
\item the selected point is an
\((\omega(k)+2\eta(k))\)-approximate minimizer; and
\item the finite comparison theorem, with \(x\) represented by an observation and its
coverage witness, is derivable in \(\iPRA\) using only bounded search and rational
arithmetic.
\end{enumerate}
\end{theorem}

\begin{proof}
Fix admissible \(x\).  By coverage, choose a net point \(p_{k,j}\) with
\(d(x,p_{k,j})\le\delta(k)\).  The modulus gives
\(f(p_{k,j})\le f(x)+\omega(k)\).  Evaluation error gives
\[
   v_{k,j}\le f(p_{k,j})+\eta(k)
            \le f(x)+\omega(k)+\eta(k).
\]
Minimality of \(i_*\) yields \(m_k\le v_{k,j}\), and another evaluation estimate gives
\[
   f(p_k)\le m_k+\eta(k)
           \le f(x)+\omega(k)+2\eta(k),
\]
proving (i) and (iii).

For the lower bound in (ii), every net value satisfies
\(v_{k,i}\ge f(p_{k,i})-\eta(k)\ge\inf_X f-\eta(k)\), so their minimum does too.  For
the upper bound, apply the preceding inequality to points whose function values
approach the infimum, or equivalently take the infimum over \(x\) in the bound for
\(m_k\).

For (iv), avoid quantifying over an abstract real space.  The input supplies a finite
observation of \(x\) and a bounded index witnessing net coverage.  Booleanize the
coverage and rational inequalities, instantiate the modulus contract, and use the
finite-minimum comparison proved by induction over the net list.  The resulting object
formula says that the selected sample is no worse, up to the displayed error, than any
input carrying a coverage certificate.  Its external interpretation is (i).
\end{proof}

\subsection{No exact argmin is claimed}

An approximate minimum value is stable under small perturbations; an exact minimizing
point need not be.  A flat interval, two equal wells, or a symmetric problem may have
many minimizers.  Therefore the generic interface returns a value name and approximate
minimizer cloud, not an exact point.

If a quantitative uniqueness modulus is available, coherent minimizers can be
constructed.  One useful form is
\[
   f(x)\le \inf f+\alpha(p)
   \ \wedge\ 
   f(y)\le \inf f+\alpha(p)
   \quad\Longrightarrow\quad d(x,y)\le2^{-p}.
\]
Choose net/error parameters making the approximation error at most \(\alpha(p)\).
Then all approximate minimizers lie in a ball of radius \(2^{-p}\), providing the
diameter bound required by shrinking viability.

\subsection{Strong convexity as a uniqueness producer}

For a midpoint strongly convex function with parameter \(\lambda>0\), suppose
\[
  f\!\left(\frac{x+y}{2}\right)
  \le\frac{f(x)+f(y)}2-\frac{\lambda}{8}d(x,y)^2.
\]
Since the midpoint value is at least \(\inf f\), if both \(x\) and \(y\) are within
\(\alpha\) of the infimum, then
\[
  \frac{\lambda}{8}d(x,y)^2\le\alpha,
  \qquad d(x,y)\le\sqrt{8\alpha/\lambda}.
\]
Choosing dyadic \(\alpha\) with an explicit rational square bound yields a
primitive-recursive uniqueness modulus.  Thus strong convexity legitimately upgrades
approximate minimizer clouds to a coherent point name.

\subsection{Other compactness templates}

The same interface yields:

\begin{itemize}
\item suprema and oscillation bounds by finite max/min folds;
\item uniform bounds from pointwise estimates plus a modulus;
\item Riemann sums and cubature from finite partitions and variation moduli;
\item Lebesgue-number-style finite refinements when an explicit cover margin is given;
\item finite simplicial approximations of compact metric presentations; and
\item finite subcover conclusions when a coverage certificate is part of the input.
\end{itemize}

It does not extract a modulus or finite net from a bare qualitative assertion of
compactness.  The net producer is mathematical data.

\section{Guarded operations on represented quantities}
\label{sec:guarded}

Approximate names support algebraic operations only when error propagation is
quantified.  Operations with singularities require explicit guards.  This section
states the generic reciprocal theorem and the broader interface discipline it
illustrates.

\subsection{Regular rational names}

A convenient real presentation is a rational evaluator \(a(k)\) with a fixed
coherence estimate, for example
\[
  \abs{a(k)-a(\ell)}\le2^{-k}+2^{-\ell}.
\]
Equivalent regular-name conventions are possible.  The important point is that a
finite output demand gives a primitive-recursive input precision through an error
budget.  Addition uses the sum of errors.  Multiplication additionally needs a bound
on the magnitudes of the factors.  Polynomial evaluation uses a finite telescoping
estimate or Horner recurrence with an explicit coefficient/value bound.

\subsection{Nonzero is not apartness}

Constructively, the proposition \(x\ne0\) is negative information.  It does not supply
a rational \(\lambda>0\) with \(\abs{x}\ge\lambda\).  A uniform operation that turned
an arbitrary represented nonzero real into such a \(\lambda\) would decide information
not continuously available from a Cauchy name.  Therefore reciprocal is a guarded
operation.

\begin{definition}[Apartness witness]
An apartness witness for a represented real \(x\) is positive rational data
\(\lambda>0\) together with a certificate \(\abs{x}\ge\lambda\).
\end{definition}

\begin{theorem}[Guarded reciprocal]
\label{thm:guarded-reciprocal}
Assume \(\abs{x}\ge\lambda>0\) and a rational approximation \(q\) satisfies
\(\abs{x-q}\le\eps\le\lambda/2\).  Then
\[
  \abs{q}\ge\lambda/2
\]
and
\[
  \left|\frac1x-\frac1q\right|
  \le\frac{2\eps}{\lambda^2}.
\]
Consequently a regular name for \(x\), together with \(\lambda\), yields a
primitive-recursive regular name for \(1/x\).
\end{theorem}

\begin{proof}
The reverse triangle inequality gives
\[
  \abs{q}\ge\abs{x}-\abs{x-q}
            \ge\lambda-\eps\ge\lambda/2.
\]
Therefore
\[
 \left|\frac1x-\frac1q\right|
 =\frac{\abs{x-q}}{\abs{x}\abs{q}}
 \le\frac{\eps}{\lambda(\lambda/2)}
 =\frac{2\eps}{\lambda^2}.
\]
All quantities may be replaced by integer numerator inequalities after clearing
positive denominators.  The object proof uses only ordered-ring arithmetic and the
apartness certificate.
\end{proof}

To obtain output error at most \(2^{-p}\), it suffices to request an input approximation
with
\[
  \eps\le\min\left(\frac\lambda2,
                        \frac{\lambda^2}{2^{p+1}}\right).
\]
A dyadic lower approximation to the right-hand side gives a primitive-recursive
precision transformer.

\subsection{Totalized guarded functions}

Every primitive-recursive symbol must be total.  Define a totalized reciprocal
approximator that returns a default value when the numerical guard check fails.  Its
contract has the form
\[
  \operatorname{Apart}(x,\lambda)
  \longrightarrow
  \operatorname{RecipCorrect}(x,\lambda,\operatorname{recip}(x,\lambda,p),p).
\]
The behavior outside the guard is irrelevant.  This is preferable to pretending the
operation is partial or hiding the guard in a semantic type class.

The same pattern applies to:

\begin{itemize}
\item division by a determinant or pivot;
\item normalization of a nonzero vector;
\item logarithms on a positive interval;
\item contour integrals away from poles;
\item projectors separated by a spectral gap; and
\item local inverse functions with a derivative lower bound.
\end{itemize}

\subsection{Pole clearing and algebraic reformulation}

Sometimes the guard can be avoided by changing the observable.  Instead of proving an
identity containing \(1/x\), multiply through by a power of \(x\) and prove a
polynomial identity.  Instead of selecting a normalized eigenvector, prove a statement
about the rank-one projector after clearing a norm denominator, or about the invariant
subspace.  Monic polynomial presentations similarly avoid division by a leading
coefficient.

This is not merely a workaround.  In a weak constructive theory, polynomial and
finite-algebra certificates are often the intrinsic content of a theorem, while a
normalized representative is an additional choice requiring apartness.

\subsection{Branches and exact discrete data}

Square roots, arguments, logarithms, and Puiseux leading terms combine continuous
approximation with discrete branch data.  A branch interface must include one of:

\begin{enumerate}[label=(\roman*),leftmargin=2.4em]
\item an exact algebraic presentation that determines the branch;
\item a rational region separated from every branch boundary;
\item overlapping relational branches with a stable downstream observable; or
\item an explicit boundary assumption supplying rank, sign, or monodromy data.
\end{enumerate}

An arbitrary Cauchy name does not generally determine exact discrete information by a
uniform finite procedure.  The obstruction is addressed explicitly in
\cref{sec:barriers}; the correct response is to strengthen the presentation, not to
hide an undecidable comparison inside a selector.

\section{Exact barriers and honest unresolved inputs}
\label{sec:barriers}

A general framework is useful only if it distinguishes formalization problems from
missing mathematics and from impossible interfaces.  This section records the most
important boundaries.  None of them says that the associated theorem is unprovable in
\(\iPRA\); each says that a particular generic route cannot supply the required data.

\subsection{Cardinality does not uniformly produce a cutoff}

Suppose \(P_e(n)\) is a decidable family with at most one solution.  It is tempting to
expect a uniform algorithm that, from \(e\) and the cardinality bound, returns \(B(e)\)
such that every solution is below \(B(e)\).  This is impossible.

\begin{barrier}[No computable cardinality-to-cutoff operator]
\label{barrier:cardinality-cutoff}
There is no total computable operator which, given a code for a decidable predicate
\(P(n)\) together with a valid certificate that \(P\) has at most one solution, always
returns a bound \(B\) such that
\[
  P(n)\quad\Longrightarrow\quad n<B.
\]
In particular, no primitive-recursive operator has this property.
\end{barrier}

\begin{proof}
Let \(M_e\) be the Turing machine with code \(e\).  Define the decidable predicate
\[
  P_e(n)\;:\Longleftrightarrow\;
  M_e\text{ halts for the first time after exactly }n\text{ steps}.
\]
Bounded simulation makes \(P_e\) primitive recursive in \((e,n)\).  It has at most one
solution, and the proof of that fact is uniform: a deterministic machine cannot have
two distinct first halting times.

Apply the hypothetical operator to obtain \(B(e)\).  Simulate \(M_e\) for \(B(e)\)
steps.  If it halts, answer yes.  If it does not halt, it can never halt later, because
a later first halting time \(n\ge B(e)\) would be a solution of \(P_e\) outside the
promised cutoff.  This decides the halting problem, a contradiction.
\end{proof}

Thus a finiteness or cardinality theorem is not enough for complete enumeration.  A
height bound, search cutoff, or complete-list certificate is genuinely additional
data.  This is the precise issue in generic forms of Siegel, Faltings, and
Mordell--Weil enumeration.

\subsection{Exact zero tests from approximate names}

\begin{barrier}[No finite-use exact zero test]
\label{barrier:zero-test}
There is no extensional finite-use procedure that, from an arbitrary regular Cauchy
name of a real number known only to be dyadic rational, always returns the exact bit
\([x=0]\).
\end{barrier}

\begin{proof}
Run the supposed procedure on a regular name for zero.  Because the output is a finite
computation, it reads only rows through some horizon \(H\).  Choose a positive dyadic
\(\delta=2^{-M}\) with \(M\) so large that a regular name for \(\delta\) may output
zero on every row through \(H\) and switch to sufficiently accurate approximations to
\(\delta\) later.  For example, take \(M>H+2\); the discrepancy from zero at the early
rows is below their allowed error.  The two names are indistinguishable to the
procedure, so it returns the same bit, although one value is zero and the other is not.
Contradiction.
\end{proof}

Promises about the abstract value do not repair the information interface.  Saying
``the represented number happens to be algebraic'' or ``dyadic'' is weaker than
providing a finite exact presentation from which equality is decidable.  Exact Puiseux
leading exponents, multiplicities, branch ranks, and similar discrete outputs can
encode such a zero test.  Therefore an unconditional exact-output theorem for arbitrary
Cauchy coefficients is not available.  One must require exact algebraic splitting
data or return a relational approximate object.

\subsection{No universal primitive-recursive evaluator}

The diagonal proof in \cref{barrier:no-universal-pr} blocks another common shortcut.
A code for an arbitrary primitive-recursive function is not automatically executable
by a total primitive-recursive universal semantics.  Static specialization, finite
query tables, or a restricted fuel-bounded interpreter are the admissible alternatives.

\subsection{Metastability does not imply a tail rate}

Finite-window stability is compatible with arbitrarily late future oscillations.  A
uniform conversion from every metastability bound to a Cauchy modulus would therefore
supply information not present in the premise.  In applications, a metastable theorem
is often exactly what a finite consumer needs, and it should be retained as such.  A
limit name requires a separate tail modulus, summable motion estimate, monotonicity
argument, or other global input.

\subsection{Negative nonzeroness does not supply positive separation}

The guarded reciprocal theorem shows what division needs: a rational lower bound
\(\lambda\).  The proposition \(x\ne0\) alone gives no algorithm for finding
\(\lambda\) from an arbitrary Cauchy name.  This is another discontinuity at zero.
Accordingly, every inverse, normalization, pivot, or branch-isolation interface lists
its apartness data explicitly.

\subsection{Qualitative compactness does not supply a net}

A theorem that a space is compact in a strong ambient theory may have no
primitive-recursive finite-net producer with a modulus usable in \(\iPRA\).  The
\(\FRNF\) compactness interface is conditional on such data.  For boxes, simplices,
finite complexes, rational polytopes, and many explicitly presented compact spaces,
the data are easy to construct.  For an arbitrary quotient or function space, they may
be the main mathematical theorem.

\subsection{Deep analytic and arithmetic producers}

Finite contour integration, exact rational linear algebra, and bounded enumeration can
formalize the consequences of a zero-free region, exceptional-zero separation, height
bound, or covering list.  They do not manufacture those inputs.  In particular:

\begin{itemize}
\item stronger prime-number-theorem rates require correspondingly strong zero-free
regions and growth estimates;
\item Siegel--Walfisz and Bombieri--Vinogradov uniformity require substantive control of
exceptional characters/zeros, not merely better bookkeeping;
\item general Chebotarev needs effective field algorithms and analytic estimates with
field dependence tracked;
\item general Siegel and Faltings enumeration needs an effective height producer; and
\item general central-limit and stable-law statements need tail, smoothing,
quantization, or classification inputs beyond finite-law arithmetic.
\end{itemize}

\subsection{Effectivity obstacles are not independence proofs}

Failure of a generic producer does not prove that no \(\iPRA\) proof of a particular
closed theorem exists.  A theorem may have a special argument, a nonuniform constant,
or a certificate discovered by external mathematics and then verified in \(\iPRA\).
The barriers above justify exact claim wording:

\begin{quote}
``The finite verifier and conditional implication are formalized; the required
producer is not supplied by this framework.''
\end{quote}

They do not justify the stronger wording ``the theorem is unprovable in \(\iPRA\)''
without a separate proof-theoretic argument.

\section{Domain interfaces and reusable proof recipes}
\label{sec:domains}

The framework is deliberately domain-neutral, but its value is clearest when the
abstract nodes are translated into concrete mathematical recipes.  This section gives
such recipes and identifies the remaining input for each area.

\subsection{Finite algebra and combinatorics}

Exact finite algebra is the easiest case.  Encode tuples, lists, matrices, polynomials,
permutations, graphs, and finite sets by natural numbers.  Operations are primitive
recursive, and identities are equations or bounded formulas.  The standard recipe is:

\begin{enumerate}[label=\arabic*.,leftmargin=2.2em]
\item choose a canonical code or a proof-relevant equality certificate;
\item define operations by fixed primitive recursion;
\item prove defining equations and congruence in \(\Prf\,\PRAAx\);
\item Booleanize bounded side conditions;
\item prove the main invariant by QF induction on size, degree, or a loop clock; and
\item use the proof schedule to compose algebraic subroutines.
\end{enumerate}

Determinants by Laplace expansion, adjugates, bounded sums/products, Bell recurrences,
finite group actions, and combinatorial bijections fit directly.  Symbolic polynomial
identities may be proved in a host ring and reflected through a checked object-level
ring lemma, as the repository's \code{ipra\_ring} tactic already demonstrates.

\subsection{Exact coordinate geometry and general relativity identities}

Incidence and metric identities over rational or integer coordinates reduce to
polynomial equalities and guarded divisions.  Degenerate configurations are handled by
explicit nonvanishing hypotheses or by clearing denominators.  For symbolic spacetime
metrics, one fixes a coordinate chart and rational/algebraic parameter coding, computes
inverse matrices, Christoffel symbols, curvature, and field tensors by finite sums, and
proves the resulting polynomial/rational identities under denominator guards.

The framework proves local coordinate identities.  It does not infer global causal or
manifold statements unless chart coverage, transition maps, and topological
certificates are supplied as separate nodes.

\subsection{Algebraic numbers, roots, and spectral data}

A robust root interface returns a finite cloud of rational disks or boxes with:

\begin{itemize}
\item residual bounds;
\item a coverage/count certificate;
\item overlap or continuation edges across precision;
\item symmetric-product or coefficient reconstruction estimates; and
\item separation data only when exact individual labels are requested.
\end{itemize}

Multiple roots and eigenvalue collisions are treated relationally.  Stable outputs
include the unordered multiset, elementary symmetric functions, cluster sums,
characteristic polynomials, and spectral projectors onto isolated clusters.  An
ordered eigenvector or individual branch requires a gap and normalization guard.

Variable degree is handled by one clocked recursion on degree after a fixed coefficient
evaluator has been statically specialized.  Backward dual-horizon planning determines
all coefficient rows needed for repeated deflation and residual verification.

\subsection{Real analysis and calculus}

A function presentation should contain the quantitative data familiar from constructive analysis \cite{bishopbridges1985,troelstravanDalen1988}:

\begin{itemize}
\item a rational sampler or regular-name evaluator;
\item a domain certificate;
\item a continuity, derivative, variation, or holomorphy modulus appropriate to the
theorem;
\item magnitude bounds needed for products and remainders; and
\item finite algorithms for partitions, sums, or polynomial approximants.
\end{itemize}

The fundamental theorem of calculus, Taylor bounds, Euler--Maclaurin, and special-
function identities are then stated as rational inequalities at requested precision.
One computes a partition size or truncation order from the supplied moduli, verifies a
finite sum identity, and bounds the remainder.  No induction is performed on the
predicate ``this is a Cauchy name''; the final precision is frozen and only the finite
rows selected by the justification horizon occur in the invariant.

Qualitative differentiability or holomorphy is not silently converted to a uniform
modulus.  For specially constructed functions, a modulus may be derived from explicit
convexity, coefficient, or derivative bounds and then becomes another compiled node.

\subsection{Complex analysis and contour methods}

Contours are finite polygonal or rational parametrized paths with derivative and
clearance bounds.  Analytic functions are represented on a certified neighborhood by
samplers and uniform error/modulus data.  Cauchy and residue arguments become finite
quadrature plus explicit deformation or local-series estimates.  Poles are guarded by
separation from the contour.  Winding numbers can be returned as exact integers when a
positive distance from zero ensures stability under approximation.

An abstract identity theorem or unrestricted contour deformation is not needed for a
fixed application: backward demand propagation expands the finite tree of local
comparisons and approximation estimates needed by the final contour calculation.

\subsection{Analytic number theory}

Dirichlet polynomials, finite Euler products, summatory functions, and truncated
integrals are finite arithmetic.  The difficult inputs are explicit analytic estimates.
An \(\FRNF\) theorem should state them as contracts with all dependencies visible:
modulus, conductor, field discriminant, strip, height, or zero-free-region parameters.
The finite compiler then:

\begin{enumerate}[label=(\roman*),leftmargin=2.4em]
\item selects truncation and contour parameters;
\item allocates precision backward through the finite expression graph;
\item verifies rational/complex approximations and denominator guards;
\item carries error numerators by QF induction; and
\item emits the final explicit inequality.
\end{enumerate}

This architecture can strengthen reliability and reuse without pretending to improve
the underlying zero-free region.

\subsection{Diophantine geometry and complete enumeration}

Given a family \(X_p\) and a primitive-recursive height bound \(H(p)\), enumerate every
rational/integral code of height at most \(H(p)\), test the defining equations, and
return the filtered list.  The object proof establishes soundness by direct checking
and completeness by the height theorem.  Descent and local solubility certificates may
reduce the search or prove the bound in special families.

Without \(H\), a finite verifier remains useful but the complete-list theorem is
conditional.  A cardinality theorem alone cannot fill the gap by
\cref{barrier:cardinality-cutoff}.

\subsection{Probability}

Finite rational probability laws are exact tables.  Expectations, moments,
convolutions, characteristic functions at rational points, and finite-dimensional
distribution comparisons are bounded sums.  Quantitative limit theorems require
explicit rates such as Berry--Esseen constants, truncation/tail bounds, Lindeberg
moduli, or smoothing inequalities.

For an infinite or real-valued law, the interface contains a primitive-recursive
quantization at precision \(k\), a tail/error bound connecting it to the original law,
and the finite-law theorem.  The proof schedule composes these.  It does not quantify
inside first-order arithmetic over arbitrary measures as higher-type objects.

\subsection{Topology}

A finite simplicial complex is exact combinatorial data, so chains, boundaries, cycles,
homology certificates, and simplicial maps are finite linear algebra.  A represented
space is observed through a tower of finite complexes.  A map returns a finite carrier
cloud rather than a forced vertex at ties.  If every selection from the carrier is
simplicial and any two selections are contiguous, an explicit prism operator gives a
chain homotopy and makes the induced homology map independent of choices.

Jordan-type and dimension arguments then reduce to finite grid parity, Sperner labels,
or chain calculations plus supplied sampling and inverse-modulus data.  The observation
tower unifies the presentation; it does not generate missing continuity moduli.

\subsection{Category theory and algebraic geometry}

For fixed finite presentations, categorical constructions are finite data:
objects/morphisms have codes, composition is a table or fixed program, and diagrams are
bounded lists.  Universal properties are represented by explicit witness maps and
uniqueness derivations.  Yoneda components, finite limits/colimits, adjunction units,
and finite Kan extensions can therefore be expressed in \(\iPRA\) for specified
families.

Quotients should be certified setoids, not opaque host quotients.  Equality in a
quotient is witnessed by a finite certificate---for example an ideal combination,
localization killing exponent, or rewrite normalization.  Affine gluing uses finite
covers, overlap maps, and cocycle certificates.  Scheme-theoretic assertions requiring
unbounded families or non-effective existence remain outside the finite presentation
unless additional producers are supplied.

\subsection{A common audit question}

Across all areas, the decisive question is:

\begin{quote}
For the final numerical challenge \(d\), can every object, query, branch, and proof
obligation actually used by the argument be placed in one primitive-recursively bounded
finite graph, with local correctness derivations and explicit external producers?
\end{quote}

If yes, the proof is a candidate for \(\FRNF\).  If no, the failure identifies either a
missing quantitative theorem or an inappropriate representation.

\section{The Finitary Realizer Normal Form compilation theorem}
\label{sec:compilation-theorem}

We now collect the preceding constructions into one theorem schema.  The statement is
metatheoretic: it takes fixed object syntax and local embedded derivations and returns a
new embedded derivation.  It does not assert that \(\iPRA\) internally quantifies over
all such graphs.

\subsection{FRNF graph}

\begin{definition}[Certified FRNF graph]
A certified \(\FRNF\) graph consists of the following finite data.

\begin{enumerate}[label=\textnormal{G\arabic*.},leftmargin=3.4em]
\item A finite acyclic directed graph with a designated final node and fixed
 topological order.
\item For each input, a finite-observation presentation: exact data, regular names,
 finite clouds, finite nets, or certified setoids.
\item For every edge, primitive-recursive computation and justification horizon
 transformers with an embedded proof that read demand is bounded by justification
 demand.
\item For every node, fixed primitive-recursive syntax and an object contract formula.
\item Each computational node is one of: exact transducer, bounded search, clocked
 loop, finite relational viability pass, stable observer, gap collapse, boundary node,
 or finalizer.  Compound domain routines are expanded into these forms.
\item Every clocked loop has a bounded/QF invariant and embedded base, step, and finish
 derivations.
\item Every bounded formula occurring in a local invariant has a characteristic term
 and the embedded Booleanization derivations of \cref{thm:booleanization}.
\item Every node has an embedded local derivation of its contract from its declared
 predecessor contracts and permanent boundary formulas.
\item Every relational output consumed deterministically is either internal to a
 finite viability computation or followed by an embedded stability/gap derivation.
\item Every non-finite mathematical input---modulus, height bound, separation, finite
 net, zero-free estimate, covering list---is either produced by a preceding certified
 node or appears explicitly among the boundary formulas.
\end{enumerate}
\end{definition}

The definition is intentionally syntactic.  Host-semantic theorems may accompany the
graph, but they are not among the premises needed to construct the object proof.

\subsection{Compilation theorem}

\begin{theorem}[FRNF compilation]
\label{thm:frnf-compilation}
Let \(G\) be a certified \(\FRNF\) graph and let \(d\) be a final demand term.  There is
a metatheoretically effective construction that produces:

\begin{enumerate}[label=(\roman*),leftmargin=2.4em]
\item primitive-recursive terms coding the least dual-horizon demand closure of
\((G,d)\);
\item finite input-table and trace terms sufficient for every computation and proof in
\(G\);
\item an embedded derivation
\[
   \Gamma_G\vdash_{\PRAAx}\operatorname{Goal}_G(d),
\]
where \(\Gamma_G\) is exactly the list of undischarged boundary formulas; and
\item by the existing transfer theorem, a derivation of the translated goal in the
native scalar core.
\end{enumerate}
No induction beyond quantifier-free induction is introduced by the construction.
\end{theorem}

\begin{proof}
We describe the compiler and verify each stage.

\paragraph{1. Static specialization.}
Replace every externally fixed routine by concrete primitive-recursive syntax and its
defining equations.  If a routine was supplied as a finite table or as a code in a
fixed restricted language, instantiate the corresponding table lookup or bounded
interpreter.  No total universal evaluator is introduced.

\paragraph{2. Reverse dual-demand pass.}
Apply \cref{thm:demand-closure} to the finite graph at final demand \(d\).  This yields
for each node/input the least read and justification bounds.  The pass itself is a
bounded primitive recursion on the encoded reverse topological list; its correctness
has a QF invariant.

\paragraph{3. Finite materialization.}
Construct codes for every input row, validity instance, approximation, candidate, and
edge relation at indices within the justification closure.  Exact inputs are constant
tables.  For regular names, the global name hypotheses are instantiated only at the
finitely many requested indices.  For relational inputs, materialize the bounded
candidate tables.

\paragraph{4. Booleanization.}
Apply \cref{thm:booleanization} to every bounded validity, range, table, and edge
condition.  Local state predicates become equations \(\chi=1\), and finite universal
conditions become bounded-product equations.  Thus all loop invariants are
quantifier-free.

\paragraph{5. Local computational nodes.}
Exact transducers use defining equations.  Bounded searches use
\cref{thm:bounded-search}.  Majorized descents use
\cref{thm:majorized-descent}.  Finite relational nodes use
\cref{thm:backward-viability}.  Each adaptive routine has been defunctionalized into a
clocked transition system, so \cref{thm:trace-internalization} constructs its embedded
contract from the supplied base, step, and terminal derivations.

\paragraph{6. Choice erasure.}
When a downstream observable consumes relational candidates, use the supplied embedded
stability theorem and \cref{thm:stable-observer}.  If an exact discrete label is
requested, use the explicit gap data and \cref{thm:gap-collapse}.  No arbitrary exact
branch choice is generated.

\paragraph{7. Topological proof composition.}
Process graph nodes in topological order.  By assumption, each local contract is
derivable from declared predecessors and \(\Gamma_G\).  Apply
\cref{thm:proof-schedule} at every node.  Induction on the finite list yields a
derivation of the final goal from \(\Gamma_G\).

\paragraph{8. Native export.}
Apply \code{pra\_iff\_native} to the completed ordinary-iPRA derivation.  This stage
changes presentation but not mathematical assumptions.

Every induction introduced above is on a numerical clock/table index with a
Booleanized state equation.  Hence it is quantifier-free induction.  All other logical
steps are ordinary intuitionistic deduction and equality reasoning.  The only
undischarged hypotheses are the boundary formulas because all generated component
contracts have been cut away.
\end{proof}

\subsection{Relative expressiveness}

The theorem is not a completeness theorem for all constructive mathematics.  It is a
normal-form theorem for proofs already equipped with primitive-recursive finite
realizers and quantitative interfaces.  Within that class, the node vocabulary is
broad: every fixed primitive recursion is a clocked loop; every bounded existential is
least search; every finite nondeterministic process is a layered relation; every finite
DAG composes by the schedule.  Thus the remaining expressiveness questions are mainly
about whether the mathematics supplies finite bounds and stable representations.

\subsection{Conservativity of the compiler}

The compiler adds no axioms to \(\iPRA\).  Its generic Lean definitions live in the
metatheory and build terms of the inductive relation \(\Prf\,\PRAAx\).  Trust rests on
the Lean kernel, the definition of \(\Prf\), and the audited local constructors.  A
future executable compiler may be large or use sophisticated metaprogramming without
changing the object-theory theorem, provided the generated proof term checks.

\subsection{The main engineering gap}

The logical schedule and QF loop combinators are short.  The substantial remaining
implementation is a reusable object arithmetic layer for:

\begin{itemize}
\item generic bounded-formula reification and substitution;
\item encoded arrays and finite folds;
\item graph demand tables;
\item viability masks and path extraction; and
\item rational error budgets and metric inequalities.
\end{itemize}

The checked \code{IPRADerivations} library already contains most low-level ingredients.
The task is integration and abstraction, not a new foundational theory.

\section{Resolution map for the reported difficulties}
\label{sec:resolution-map}

The repository audit revealed a stable list of recurring obstacles.  The table below
records what the framework resolves, what it resolves only under explicit data, and
what remains substantive mathematics.  ``Resolved'' means that the generic reduction
is proved; it does not mean that every application packet has been migrated to a
kernel-checked object proof.

\small
\begin{longtable}{@{}p{0.16\textwidth}p{0.23\textwidth}p{0.32\textwidth}p{0.20\textwidth}@{}}
\toprule
Obstacle & Invalid or incomplete move & FRNF repair & Status / remaining input \\
\midrule
\endfirsthead
\toprule
Obstacle & Invalid or incomplete move & FRNF repair & Status / remaining input \\
\midrule
\endhead

Unbounded precision induction & Put a predicate quantifying over all rows of a Cauchy
name inside induction. & Freeze the final demand, compute the full dual-horizon closure,
materialize its finite table, and Booleanize all table contracts. & Generic reduction
resolved; object table library still to implement. \\
\addlinespace

Adaptive oracle use & Let future queries depend on higher-type answers without a fixed
finite domain. & Defunctionalize to a bounded transition system; compute worst-case
read and proof horizons; verify the actual trace. & Resolved when an explicit clock and
query majorant are supplied. \\
\addlinespace

Computation versus proof horizon & Track only rows read by the executable routine. &
Carry separate read and justification functions and propagate both backward. & Resolved
architecturally and mathematically. \\
\addlinespace

Exact boundary or chart choice & Decide sign, equality, or chart membership for an
arbitrary represented real. & Use overlapping branches or finite candidate clouds;
collapse only from a positive gap. & Conditional resolution; gap required for exact
label. \\
\addlinespace

Nonzero to inverse & Treat \(x\ne0\) as a rational lower bound. & Guard inverse by
\(\abs{x}\ge\lambda>0\), or clear denominators and prove a polynomial identity. &
Requires apartness input unless reformulated. \\
\addlinespace

Abstract compactness & Invoke compactness, attainment, or a finite subcover without
quantitative data. & Certified finite nets, coverage witnesses, and moduli; return
approximate extrema. & Conditional resolution; net/modulus producer required. \\
\addlinespace

Root and eigenvector labels & Globally order branches and vectors through
multiple roots or eigenspaces. & Return clusters, symmetric observables, projective
objects, or isolated spectral subspaces. & Resolved for invariant outputs; individual
label needs separation. \\
\addlinespace

Metastability to convergence rate & Infer an eventual tail modulus from one stable
finite window. & Keep the metastable contract; derive a Cauchy name only from a tail or
summable-motion modulus. & Exact boundary retained. \\
\addlinespace

Cardinality to cutoff & Stop after a number of found solutions equal to the claimed
cardinality. & Require a height/cutoff producer or a complete-list certificate. & No
generic operator; impossible by \cref{barrier:cardinality-cutoff}. \\
\addlinespace

Host Lean theorem to iPRA & Call a theorem about \(\N\), \(\Q\), or \(\mathbb R\) an
object proof. & Formula firewall: concrete PR syntax, object formula, local
\(\Prf\,\PRAAx\) constructor, proof schedule. & Engineering target; logical kernel
combinators supplied. \\
\addlinespace

Universal PR evaluator & Execute an arbitrary primitive-recursive program code inside
one PR function. & Statically specialize fixed evaluators, use finite tables, or a
restricted fuel-bounded interpreter. & Exact impossibility by diagonalization. \\
\addlinespace

Arbitrary Cauchy Puiseux data & Return exact leading exponent/coefficient from an
approximate coefficient name. & Require finite exact algebraic splitting/rank data, or
return residual clusters/regular branches. & Exact-output interface obstructed by the
zero-test barrier. \\
\addlinespace

Qualitative differentiability or holomorphy & Extract a uniform derivative or analytic
modulus from a bare qualitative assertion. & Supply a modulus or derive one from the
special presentation; then use finite partitions/series and explicit remainders. &
Conditional resolution. \\
\addlinespace

General Siegel/Faltings & Treat finite verification as a source of a height bound. &
Compile ``height bound implies complete list''; keep the bound producer visible. &
Finite consequence resolved; general producer remains deep mathematics. \\
\addlinespace

Strong PNT/uniform prime distribution & Expect demand bookkeeping to improve a weak
zero-free region or remove exceptional zeros. & Express every analytic estimate as a
boundary contract and preserve parameter dependence. & Formalization improved; stronger
analytic theorem remains. \\
\addlinespace

General probability laws & Quantify over arbitrary measures/functions and use weak
compactness. & Finite rational laws plus explicit quantization, tail, Lindeberg, and
smoothing moduli. & Conditional resolution; domain estimates required. \\
\addlinespace

Quotient equality & Rely on opaque host quotients and extensional equality. & Canonical
normal forms or certified setoids with finite equality witnesses. & Resolved as an
interface discipline. \\
\addlinespace

Infinite branch selection & Apply K\"onig's lemma to compatible finite levels. & Finite
backward viability for each horizon; coherence from shrinking diameter and motion
budgets. & Resolved for quantitative represented outputs. \\
\bottomrule
\end{longtable}
\normalsize

\subsection{Resolved does not mean automatic}

For example, the dual-horizon theorem gives a correct composition rule, but an
application still has to derive its local horizons.  The finite-net theorem gives an
approximate minimum once a net and modulus exist, but it does not construct those for
an arbitrary compact object.  The proof schedule builds a final derivation once local
\(\Prf\,\PRAAx\) proofs are supplied, but it does not synthesize domain lemmas from a
host proof.  The framework turns vague gaps into finite, typed obligations.

\subsection{Highest-value migration targets}

The audit suggests the following priority order.

\begin{enumerate}[label=\arabic*.,leftmargin=2.2em]
\item Complete the generic bounded-formula and encoded-table object library.
\item Formalize reverse demand closure and the two-pass viability algorithm as actual
\(\Prf\,\PRAAx\) theorems.
\item Migrate one representative analytic packet whose moduli are already explicit,
so that dual horizons are exercised end to end.
\item Migrate one relational application with collisions, such as stable root clusters
or simplicial carriers.
\item Express a conditional height-to-list theorem and connect a special-family height
producer.
\item Only then generalize the code generator to larger packet families.
\end{enumerate}

These tasks maximize reuse and directly attack the semantic-Lean-to-object-iPRA gap
identified throughout the corpus.

\section{Lean formalization plan}
\label{sec:lean-plan}

This section gives a concrete path from the mathematical framework to a complete Lean
formalization in the repository.  The target environment is Lean 4.32.0 \cite{lean4} with the
matching mathlib revision, as pinned by the current ProveIt root.  The intended library
location is
\[
\code{Logic/PrimitiveRecursiveArithmetic/Derivations/Lean/IPRADerivations/FRNF/}.
\]

\subsection{Reuse before abstraction}

The implementation should reuse the checked modules rather than create parallel
arithmetic:

\begin{center}
\begin{tabularx}{\textwidth}{@{}p{0.34\textwidth}X@{}}
\toprule
Existing module & Reuse \\
\midrule
\code{Core/Basic} & \code{inst}, \code{ind}, defining equations, congruence. \\
\code{Core/Semiring}, \code{Core/Order}, \code{Core/Monus} & Boolean arithmetic,
max/min, inequalities, error budgets. \\
\code{Core/Pairing} & Pair codes, projections, equality test, simultaneous recursion. \\
\code{Core/Search} & Least bounded search and minimality. \\
\code{Core/BoundedSums}, \code{Core/MultiSums} & Booleanized bounded quantifiers and
finite folds. \\
\code{Core/Definitions}, \code{Core/Sums} & Generated primitive-recursive recurrences. \\
\code{Core/Descent}, \code{Core/DescentEx} & Generic finite/infinite descent endpoints. \\
\code{Core/Division}, \code{Core/Gcd}, \code{Core/Bezout} & Majorized Euclidean examples
and denominator arithmetic. \\
\code{NativeBiInterpretation} & Final \code{pra_iff_native} export. \\
\bottomrule
\end{tabularx}
\end{center}

\subsection{Proposed module sequence}

\paragraph{1. \code{FRNF/BoundedFormula.lean}.}
Define the bounded formula syntax independently of host propositions:

\begin{lstlisting}
inductive BForm : Nat -> Type
  | eq    (s t : PTm n)
  | bot
  | and   (A B : BForm n)
  | or    (A B : BForm n)
  | imp   (A B : BForm n)
  | allLt (bound : PTm n) (body : BForm (n+1))
  | exLt  (bound : PTm n) (body : BForm (n+1))
\end{lstlisting}

Implement \code{toFm}, generated characteristic symbols, and substitution.  Prove
simultaneously:

\begin{lstlisting}
theorem char_boolean (A : BForm n) :
  [] |-[(PRAAx)] (charTm A =' 0) \/ (charTm A =' 1)

theorem to_char (A : BForm n) :
  [A.toFm] |-[(PRAAx)] charTm A =' 1

theorem of_char (A : BForm n) :
  [charTm A =' 1] |-[(PRAAx)] A.toFm
\end{lstlisting}

The bounded quantifier cases should invoke generic bounded sums/products.  Prove
\code{char_subst} early; otherwise every later generated invariant becomes painful.

\paragraph{2. \code{FRNF/Table.lean}.}
Choose one list/array coding based on pairing and simultaneous recursion.  Define
\code{len}, \code{get}, \code{set}, \code{append}, \code{foldMax}, and a bounded map.
Prove object equations for in-range lookup, update, and fold invariants.  Avoid a
universal function argument: instantiate folds with fixed bodies or a fixed finite
opcode family.

\paragraph{3. \code{FRNF/Horizon.lean}.}
Represent each concrete horizon by a \code{PR} symbol.  Define prefix maximum and prove
monotonicity, domination, and dual-horizon composition.  A useful object record is a
metatheoretic structure whose fields include the symbols and the embedded derivations:

\begin{lstlisting}
structure DualHorizonCert (r j : PR 1) where
  read_le_justify : [] |-[(PRAAx)]
    allOpen (app r ![var 0] <=' app j ![var 0])
  read_mono : ...
  justify_mono : ...
\end{lstlisting}

The structure itself is host data, but every field used by the compiler is an embedded
proof.

\paragraph{4. \code{FRNF/Demand.lean}.}
For a fixed generated graph, first formalize a specialized demand pass.  Then factor a
generic encoded-graph theorem if worthwhile.  The main endpoint is a QF loop proving
that the reverse array after \(m\) steps is closed and least.  Separate read and proof
arrays in the state.

\paragraph{5. \code{FRNF/Trace.lean}.}
Move the source-level \code{QFLoop} and proof-schedule combinators from this archive
into the checked library.  Add a record for a concrete transition with:
\code{state}, \code{init}, \code{step}, \code{fuel}, \code{inv}, and embedded base/step
proofs.  The compiler theorem should return \code{Prf PRAAx}, not a host predicate.

\paragraph{6. \code{FRNF/Viability.lean}.}
Encode level offsets, candidate counts, validity bits, edge bits, and viability masks.
Implement the backward dynamic program and forward least path.  Prove:

\begin{lstlisting}
theorem viable_iff_suffix_path : ...
theorem selected_path_valid : ...
theorem selected_path_compatible : ...
\end{lstlisting}

The first theorem is most easily proved by QF induction on the reversed level index;
existence of a suffix path is represented by a bounded Boolean characteristic, not an
unbounded existential in the invariant.

\paragraph{7. \code{FRNF/Stability.lean}.}
Define finite diameter, stable observers, gap collapse, and the shrinking-viability
one-step estimate.  Initially specialize to rational observables encoded by signed
integer pairs.  General metric presentations can be layered later.

\paragraph{8. \code{FRNF/FiniteNet.lean}.}
Define a finite rational net interface and exact minimum fold.  Prove the finite
comparison contract from explicit coverage and modulus formulas.  Keep the semantic
real infimum theorem in a separate host module; the object endpoint should quantify
only over coded inputs and witnesses.

\paragraph{9. \code{FRNF/Guarded.lean}.}
Formalize rational absolute values, positive denominators, and the cleared-denominator
version of \cref{thm:guarded-reciprocal}.  Export reusable precision transformers for
reciprocal, quotient, and normalized vectors.

\paragraph{10. \code{FRNF/Compiler.lean}.}
Define the proof schedule over formulas and a generated node descriptor.  A fully
generic reflection layer is optional.  It is sufficient that a metaprogram produces
concrete Lean definitions and applies \code{Schedule.extend}.  Add final native export.

\subsection{Formula generation versus opaque declarations}

The roadmap file in this package uses \code{opaque} formulas to make endpoints
unambiguous.  The completed implementation must replace every opaque endpoint with an
explicit \(\Fm\) expression assembled from the actual PR symbols.  Otherwise a proof
of the opaque formula would not establish the intended theorem.

A useful discipline is to define each object formula immediately next to a human
pretty-printer and a host evaluation theorem.  The pretty-printer assists review; the
host theorem assists debugging; only the explicit object formula and its embedded
proof establish the \(\iPRA\) result.

\subsection{Suggested first end-to-end milestone}

The smallest convincing milestone is not the full generic graph compiler.  It is a
single integrated theorem:

\begin{quote}
Given a fixed two-stage transducer with distinct computation and justification
horizons, compile the reverse demand bounds, materialize a finite table, verify one
clocked loop, and return a closed \(\Prf\,\PRAAx\) derivation; then export it through
\code{pra_iff_native}.
\end{quote}

This exercises every architectural boundary while keeping domain arithmetic trivial.
The second milestone should add a finite relational cloud and stable observer.

\subsection{Build and audit procedure}

A result should be labeled object-checked only after the following commands succeed in
the pinned repository environment:

\begin{lstlisting}[language=bash]
lake build IPRADerivations
lake env lean Logic/PrimitiveRecursiveArithmetic/Derivations/Lean/\
  IPRADerivations/FRNF/Audit.lean
\end{lstlisting}

The audit module should contain \code{#print axioms} for every exported theorem and a
machine-readable list of declarations containing \code{sorryAx}.  The acceptable
metatheoretic axioms should match the established repository policy; no result with
\code{sorryAx} is object-checked.

\subsection{What the accompanying Lean files establish}

\code{lean/FRNF/Kernel.lean} contains the logical schedule, QF-loop closure, stable
observer, boundary discharge, gap collapse, and native export using the inspected
repository API.  It contains no \code{axiom} declaration and no \code{sorry}, but it
was not compiled in this environment.  \code{lean/FRNF/ObjectTargets.lean} deliberately
contains visible \code{sorry} markers for Booleanization, demand closure, viability,
stable clouds, finite nets, and guarded reciprocals.  These are a roadmap, not checked
results.  This status split should be preserved during migration.

\section{Worked examples}
\label{sec:examples}

The examples below show how diverse arguments enter the same normal form.  They are
chosen to exercise exact descent, approximate compactness, relational ambiguity, and a
conditional deep theorem.

\subsection{Euclidean algorithm as a majorized trace}

Input: naturals \(a,b\).  State:
\[
  s_t=(x_t,y_t),\qquad s_0=(a,b).
\]
Transition:
\[
  (x,y)\mapsto
  \begin{cases}
    (x,0),&y=0,\\
    (y,x\bmod y),&y>0.
  \end{cases}
\]
Rank: \(\rank(x,y)=y\).  Clock: \(b+1\).  Invariant:

\begin{enumerate}[label=(\roman*),leftmargin=2.4em]
\item \(\gcd(x_t,y_t)=\gcd(a,b)\);
\item if \(y_t>0\), then \(x_t\bmod y_t<y_t\);
\item once \(y_t=0\), the state is frozen.
\end{enumerate}

The equality of gcds is expressed through the checked divisibility specification or
through preservation of common divisors.  The rank theorem gives termination by the
clock.  The terminal state satisfies \(x_t=\gcd(a,b)\).  Adding coefficient fields and
a parity bit yields the extended Euclidean algorithm and B\'ezout identity.  The
reference test on \((1071,462)\) reaches \(21\) and confirms that the trace length is
bounded by the initial second component plus the frozen step.

This example illustrates that a growing proof state need not have a decreasing total
size.  Only one certified rank field controls termination.

\subsection{Approximate minimum on a rational interval}

Let \(f:[-1,1]\to\mathbb R\) be represented by a sampler with evaluation error
\(\eta(k)\) and a Lipschitz bound \(L\).  Use the dyadic net
\[
   N_k=\{-1+i2^{-k}:0\le i\le2^{k+1}\},
\]
with radius \(\delta(k)=2^{-k}\) and modulus
\(\omega(k)=L2^{-k}\).  Evaluate every point at error \(\eta(k)\), choose the least
sample minimum, and apply \cref{thm:finite-net-min}.  The selected point \(p_k\)
satisfies
\[
   f(p_k)\le f(x)+L2^{-k}+2\eta(k)
   \quad\text{for every }x\in[-1,1].
\]

For \(f(x)=(x-1/3)^2\), the exact reference implementation on the quarter-grid chooses
\(1/4\).  An object proof need not mention the real infimum.  It proves the finite
comparison against any rational input carrying a nearest-grid witness; the semantic
corollary follows from the sampler interpretation.

If one additionally proves strong convexity, the approximate minimizer clouds have a
shrinking diameter, and \cref{thm:shrinking-viability} yields a coherent name for the
unique minimizer.  Without uniqueness, only the value and approximate cloud are
canonical.

\subsection{Polynomial roots through a collision}

Consider a monic polynomial family \(P_t(z)\) and a finite root routine that returns
rational disks covering all complex roots with multiplicity.  At a parameter where two
roots collide, no continuous exact ordering of the individual roots is expected.  The
\(\FRNF\) presentation stores:

\begin{itemize}
\item the disk cloud and multiplicity/count certificates;
\item compatibility edges between overlapping refined disks;
\item a movement bound between refinements; and
\item coefficient reconstruction bounds for the product of the disk centers.
\end{itemize}

Backward viability chooses finite compatible refinements.  The elementary symmetric
functions of any valid root multiset are stable because they reconstruct the
coefficients.  Thus the theorem about the factorization identity is exact even though
internal labels are not.  If a disk is isolated from all others by a rational gap,
\cref{thm:gap-collapse} recovers an exact branch label for that isolated cluster.

The lesson is that collision is not a failure of the root theorem.  It is a failure of
an overstrong output type.

\subsection{Variable-degree deflation with dual horizons}

Suppose a fixed coefficient evaluator gives approximations to a monic polynomial of
numerical degree \(n\).  A one-root routine at output precision \(t\) reads coefficient
rows through \(U(n,M,t)\).  Deflating by the approximate root shifts coefficient
precision and the residual proof requires comparisons through a larger bound
\(V(n,M,t)\).  Repeating the routine \(n\) times creates a changing parameter chain.

Compile it as follows.

\begin{enumerate}[label=\arabic*.,leftmargin=2.2em]
\item Define the stage state: remaining degree, normalized coefficient table, roots
already found, and error budget.
\item Compute all later verification demands backward from the final coefficient
identity.  At each stage propagate \(V\), not only \(U\).
\item Materialize the original coefficient name only through the resulting maximum.
\item Run deflation forward for exactly \(n\) stages.
\item Prove by QF induction a telescoping coefficient identity and accumulated error
bound.
\end{enumerate}

The degree is ordinary data, so this is one primitive recursion uniform in \(n\).  The
coefficient evaluator is statically fixed, respecting the no-universal-evaluator
barrier.

\subsection{Stable simplicial maps}

Let a represented map between compact spaces induce, at finite observation levels, a
carrier cloud: each source vertex maps to a nonempty set of target vertices contained
in one simplex.  Any selection gives a simplicial map.  Any two selections are
contiguous, so the explicit prism operator supplies a chain homotopy.  Therefore the
induced homology map is independent of tie-breaking.

This is a stable observer in algebraic form.  The internal vertex choices are
relational; the homology map is exact and branch-independent.  A proof requires only
finite complexes, bounded matrix operations over a finite ring, the carrier
certificate, and the prism identity.  No exact continuous selector is needed.

\subsection{Height bound implies complete Diophantine list}

Let \(F_p(\vec x)=0\) be a polynomial family and suppose a boundary formula supplies
\[
  F_p(\vec x)=0\quad\Longrightarrow\quad H(\vec x)\le B(p).
\]
Enumerate every canonical rational/integer tuple of height at most \(B(p)\), evaluate
\(F_p\) exactly, and retain the zeros.  Soundness is immediate from exact evaluation.
Completeness follows by the height formula and coverage of the enumeration.  All loops
are bounded by a primitive-recursive function of \(p\) and \(B(p)\).

The resulting theorem is
\[
  \operatorname{HeightBound}(p,B(p))
  \longrightarrow
  \operatorname{CompleteZeroList}(p,L(p)).
\]
For a special curve family, a descent theorem may provide the height bound and close
the schedule.  For general Siegel or Faltings, the finite compiler is still complete as
a conditional result, while the deep producer remains visible.

\section{Artifact validation and exact status}
\label{sec:validation}

\subsection{Reference implementation}

The file \code{reference/frnf.py} uses only the Python standard library and exact
\code{Fraction} arithmetic.  It implements:

\begin{itemize}
\item least bounded search;
\item prefix maxima and dual-horizon composition;
\item reverse demand closure on a finite DAG;
\item backward viability and least viable paths;
\item majorized frozen descent and the Euclidean example;
\item exact finite-net minimization;
\item stable-observation checking;
\item exact label agreement under a gap; and
\item the guarded reciprocal error transformer.
\end{itemize}

The eleven regression tests in \code{reference/test\_frnf.py} pass.  They include strict
separation between read and justification horizons, a shared-dependency demand graph,
a viability instance with dead ends, rejection of an empty viable path, and exact
rational error checks.  These tests validate the reference algorithms on the included
instances.  They are not proofs of the generic theorems in \(\iPRA\).

\subsection{Lean artifacts}

The Lean directory has three roles.

\begin{description}[style=nextline,leftmargin=2em]
\item[\code{FRNF/Semantic.lean}]
Small host-level records for dual horizons, static transducers, finite candidate
levels, and exact normalizers.
\item[\code{FRNF/Kernel.lean}]
Repository-aligned source for the proof schedule, QF-loop closure, stable observers,
boundary discharge, exact gap collapse, and native export.  It contains no
\code{axiom} declaration or \code{sorry} placeholder.
\item[\code{FRNF/ObjectTargets.lean}]
An explicit roadmap with the bounded-formula syntax and object theorem endpoints for
Booleanization, demand closure, viability, stable clouds, finite-net minima, and
guarded reciprocals.  It intentionally contains visible \code{sorry} placeholders.
\end{description}

Lean was not installed in the artifact runtime, so none of these new files has been
kernel checked.  In particular, absence of placeholders from \code{Kernel.lean} is not
a compilation claim.  The validation folder records every visible placeholder.

\subsection{Paper build}

The PDF is compiled from the included modular LaTeX source.  The build process runs the
Python tests, compiles the bibliography and document, checks that the PDF has extractable
text, and records hashes.  The final PDF is rendered to images and visually inspected
for clipping, missing glyphs, and broken tables.

\subsection{Claim ledger}

\code{metadata/theorem\_ledger.csv} assigns one of the following statuses to every
major theorem: fully argued, reference-tested, upstream object-checked, source-level
Lean combinator, or Lean target.  \code{metadata/obstacle\_matrix.csv} records the
resolution status and remaining input for each difficulty.  These machine-readable
files are intended to survive migration into the repository and prevent claim drift.

\subsection{Current limitations}

The package does not contain:

\begin{itemize}
\item a compiled generic \(\Prf\,\PRAAx\) Booleanization theorem;
\item an object-checked encoded table/demand/viability library;
\item a new effective height theorem for general Siegel or Faltings;
\item an exact Puiseux selector for arbitrary Cauchy coefficient names;
\item stronger analytic zero-free regions or exceptional-zero estimates; or
\item a proof that every theorem in the packet corpus fits \(\FRNF\).
\end{itemize}

What it does provide is a complete mathematical architecture, precise generic proofs,
executable finite algorithms, exact object endpoints, and a prioritized route to a
checked implementation.

\section{Conclusion and research program}
\label{sec:conclusion}

The recurring difficulty in embedding mathematics into \(\iPRA\) is not the finiteness
of the final calculation.  It is the disciplined exposure of all information that the
calculation and its proof consume.  The Finitary Realizer Normal Form makes this
exposure explicit.

A proof in \(\FRNF\) has a static finite program graph; distinct computation and
justification horizons; a reverse closure pass; finite tables rather than hidden
higher-type oracles; Booleanized bounded conditions; clocked loops with QF invariants;
relational witnesses rather than unjustified selectors; stable observers and gap
certificates; and visible mathematical producer boundaries.  The final logical step is
a proof schedule whose output is an actual embedded derivation.

This architecture resolves the broad family of formal difficulties caused by
precision-dependent induction, adaptive finite use, branch ambiguity, finite
compactness, and host/object confusion.  It also sharpens the negative conclusions.
There is no universal primitive-recursive evaluator, no generic cardinality-to-cutoff
operator, and no finite-use exact zero test for arbitrary approximate names.  These are
not reasons to abandon the program.  They tell us which representations and theorem
statements are mathematically honest.

The immediate research program is concrete.

\begin{enumerate}[label=\arabic*.,leftmargin=2.2em]
\item Complete and kernel-check bounded Booleanization and encoded finite tables in the
existing \code{IPRADerivations} library.
\item Kernel-check dual-horizon demand closure and backward viability.
\item Migrate one analytic and one relational theorem end to end, producing
\(\Prf\,\PRAAx\) and native-core endpoints.
\item Build a small code generator that emits formulas, local theorem stubs, proof
schedules, and status manifests for a static graph.
\item Develop a library of producer interfaces---finite nets, apartness, height bounds,
root clouds, simplicial carriers---with exact conditional endpoints.
\item Revisit the unresolved packets only after their missing producer is isolated in
this form; then mathematical effort can focus on the genuine estimate rather than on
repeated encoding infrastructure.
\end{enumerate}

The strongest expected payoff is cumulative.  Once a component has fixed syntax,
explicit horizons, and an embedded contract, every later theorem can reuse it by cut.
The growing body of iPRA mathematics then becomes not a collection of isolated proof
outlines, but a proof-producing library of certified finitary interfaces.

\appendix
\section{Notation and formal conventions}
\label{app:notation}

\begin{longtable}{@{}p{0.22\textwidth}p{0.70\textwidth}@{}}
\toprule
Symbol & Meaning \\
\midrule
\endfirsthead
\toprule
Symbol & Meaning \\
\midrule
\endhead
\(\iPRA\) & Intuitionistic primitive recursive arithmetic. \\
\(\PRAAx\) & The axiom family used by the embedded ordinary-iPRA calculus in
\code{ConciseIPRA}. \\
\(\Gamma\vdash_{\PRAAx}A\) & An embedded derivation, represented in Lean by
\code{Prf PRAAx Gamma A}. \\
\(\PTm n\), \(\Fm\,\PTm\,n\) & Object terms and formulas with \(n\) free variables. \\
\(\chi_A\) & Primitive-recursive characteristic term of a bounded formula \(A\). \\
\(r_F(d)\) & Largest input row read while computing output demand \(d\). \\
\(j_F(d)\) & Largest input contract row used in the proof of output demand \(d\). \\
\(W_i(a)\) & Backward viability bit for candidate \(a\) at level \(i\). \\
\(\delta_k\) & Same-level diameter bound among viable observations. \\
\(\eta_k\) & Adjacent-level motion bound along a compatibility edge. \\
\(\Gamma_G\) & Permanent boundary context of an \(\FRNF\) graph. \\
\bottomrule
\end{longtable}

Natural numbers are the object domain.  Integers, rationals, finite sequences, vectors,
matrices, and finite graphs are coded by naturals.  Rational inequalities are reduced
to signed-integer or natural-number equations after clearing positive denominators.
``Primitive recursive'' always means first-order primitive recursive over \(\N\), not
an arbitrary higher-type functional.

All generic constructions are schemas over fixed syntax.  When a program or transition
is described as a parameter in the metatheory, the object theorem uses the concrete
primitive-recursive symbol produced by static specialization.

\section{Checklist for converting a proof to FRNF}
\label{app:checklist}

\begin{enumerate}[label=\textbf{\arabic*.},leftmargin=3em]
\item \textbf{State the exact object theorem.}  Write the intended closed or
conditional first-order formula, including every modulus, bound, gap, and presentation
validity hypothesis.

\item \textbf{Classify outputs.}  Is each output exact, an approximate name, a finite
cloud, a stable observable, or an exact label under separation?  Do not use one type
for all of them.

\item \textbf{Fix program syntax.}  Replace arbitrary program parameters by static
specialization, finite tables, or a restricted interpreter with fuel.

\item \textbf{List primitive-recursive symbols.}  Give arities, recurrences, defining
equations, and required arithmetic lemmas.

\item \textbf{Derive dual horizons.}  For every input edge, compute both read and
justification bounds.  Apply prefix closure.

\item \textbf{Build the demand graph.}  Topologically order components and propagate
final demands backward.  Check shared dependencies and changing parameters.

\item \textbf{Materialize finite observations.}  Replace global name properties by
the finite row instances used within the justification closure.

\item \textbf{Booleanize bounded assertions.}  Generate characteristic terms and prove
the truth equivalences and substitution laws.

\item \textbf{Clock adaptive routines.}  Supply fuel or a decreasing rank.  Freeze
terminal states.

\item \textbf{Write QF invariants.}  Include all accumulator equations, range bounds,
validity bits, and error budgets needed by the successor proof.

\item \textbf{Keep ambiguity relational.}  Define candidate validity and compatibility;
use backward viability.  Add stability or a gap theorem before deterministic exposure.

\item \textbf{Separate producers from verifiers.}  Mark finite nets, height bounds,
apartness, zero-free regions, and covering lists as generated nodes or boundary
formulas.

\item \textbf{Construct local object proofs.}  Each theorem must conclude
\code{Prf PRAAx}; a host semantic theorem is not a substitute.

\item \textbf{Compose by schedule.}  Declare exact dependencies and apply simultaneous
cut in topological order.

\item \textbf{Audit.}  Build in the pinned toolchain, search for placeholders, print
axioms, run executable tests, and compare the final formula with the human statement.

\item \textbf{Export.}  Apply the ordinary-to-native transfer only after the ordinary
object proof is complete.
\end{enumerate}

\section{Artifact map and reproduction}
\label{app:artifacts}

The archive has the following layout.

\begin{lstlisting}
ipra_frnf_framework/
  README.md
  BUILD.md
  build.sh
  paper/
    ipra_frnf_framework.tex
    ipra_frnf_framework.pdf
    references.bib
    sections/*.tex
  lean/
    lean-toolchain
    README.md
    FRNF.lean
    FRNF/Kernel.lean
    FRNF/Semantic.lean
    FRNF/ObjectTargets.lean
  reference/
    frnf.py
    test_frnf.py
    README.md
  metadata/
    obstacle_matrix.csv
    theorem_ledger.csv
    interface.schema.json
    repository_audit.md
  validation/
    python_tests.txt
    lean_placeholders.txt
    pdfinfo.txt
    pdftotext.txt
  SHA256SUMS
\end{lstlisting}

Run the finite tests with

\begin{lstlisting}[language=bash]
cd reference
python -m unittest -v test_frnf.py
\end{lstlisting}

Build the paper with

\begin{lstlisting}[language=bash]
cd paper
latexmk -pdf -interaction=nonstopmode ipra_frnf_framework.tex
\end{lstlisting}

or run \code{./build.sh} from the archive root.  Lean compilation must be performed
after copying the proposed modules into, or otherwise linking them against, the ProveIt
repository at a compatible commit.  The archive does not vendor the repository or
mathlib.


\begin{thebibliography}{99}

\bibitem{bishopbridges1985}
E.~Bishop and D.~Bridges.
\newblock \emph{Constructive Analysis}.
\newblock Springer, 1985.

\bibitem{goedel1958}
K.~G\"odel.
\newblock \"Uber eine bisher noch nicht ben\"utzte Erweiterung des finiten
  Standpunktes.
\newblock \emph{Dialectica}, 12(3--4):280--287, 1958.
\newblock doi:10.1111/j.1746-8361.1958.tb01464.x.

\bibitem{kohlenbach2008}
U.~Kohlenbach.
\newblock \emph{Applied Proof Theory: Proof Interpretations and their Use in
  Mathematics}.
\newblock Springer Monographs in Mathematics. Springer, 2008.
\newblock doi:10.1007/978-3-540-77533-1.

\bibitem{lean4}
L.~de Moura and S.~Ullrich.
\newblock The Lean 4 theorem prover and programming language.
\newblock In \emph{Automated Deduction -- CADE 28}, pages 625--635. Springer,
  2021.
\newblock doi:10.1007/978-3-030-79876-5\_37.

\bibitem{parsons1970}
C.~Parsons.
\newblock On a number theoretic choice schema and its relation to induction.
\newblock In A.~Kino, J.~Myhill, and R.~E. Vesley, editors,
  \emph{Intuitionism and Proof Theory}. North-Holland, 1970.

\bibitem{proveit}
V.~Reshetnikov and contributors.
\newblock \emph{ProveIt: Primitive Recursive Arithmetic development}.
\newblock GitHub repository, snapshot accessed 16 November 2026; README commit
  \texttt{67f5216e0b1337b533b43064262aeff268bc56db}.
\newblock \url{https://github.com/VladimirReshetnikov/ProveIt/tree/main/Logic/PrimitiveRecursiveArithmetic}.

\bibitem{troelstravanDalen1988}
A.~S. Troelstra and D.~van Dalen.
\newblock \emph{Constructivism in Mathematics: An Introduction}, volume I.
\newblock North-Holland, 1988.

\end{thebibliography}
\end{document}
